\chapter{Uniwersum konstruowalne i modele wewnętrzne}\label{ch:constructible}
\chaptermark{Uniwersum konstruowalne}
\section{Dlaczego ograniczać zbiór potęgowy?}
Aksjomat zbioru potęgowego pozwala przejść od $X$ do wszystkich jego podzbiorów. Konstrukcja Gödla $L$ zachowuje jedynie podzbiory definiowalne nad dotychczas zbudowaną strukturą. Prowadzi to do szczególnej klasy tranzytywnej, której elementy mają kanoniczne opisy.
\begin{definicja}[Definiowalne podzbiory]\label{def:constructible-def}
Niech $M$ będzie niepustym zbiorem. Zbiór
$\operatorname{Def}(M)$ składa się ze wszystkich podzbiorów
$A\subseteq M$, dla których istnieją formuła pierwszego rzędu
$\varphi(x,z_1,\ldots,z_n)$ języka $\{\in\}$ (dopuszczamy $n=0$)
oraz parametry $p_1,\ldots,p_n\in M$ takie, że
\[
 A=\{x\in M:(M,\in\cap(M\times M))\models
       \varphi(x,p_1,\ldots,p_n)\}.
\]
Definiowalność odnosimy do \emph{struktury o uniwersum $M$},
a nie do całego uniwersum teorii mnogości.
Dla $M=\varnothing$ przyjmujemy osobno
$\operatorname{Def}(\varnothing):=\{\varnothing\}$.
Ta konwencja pozwala zdefiniować pierwszy poziom $L_1$,
bez uznawania pustego zbioru za strukturę pierwszego rzędu
(o niepustej dziedzinie w przyjętej konwencji logicznej).
\end{definicja}
\begin{lemat}[Istnienie operatora definiowalnych podzbiorów]
\label{lem:def-set}
Dla każdego zbioru $M$ istnieje w ZF zbiór
$\operatorname{Def}(M)$ i zachodzi
$\operatorname{Def}(M)\subseteq\mathcal P(M)$.
Jeśli $M\ne\varnothing$, rodzina ta zawiera $\varnothing$
oraz $M$.
\end{lemat}
\begin{proof}
Dla zbioru pustego stwierdzenie wynika bezpośrednio
z konwencji w Definicji~\ref{def:constructible-def}.
Załóżmy więc $M\ne\varnothing$. Kody formuł tworzą zbiór
liczb naturalnych, a skończone krotki parametrów z $M$
tworzą zbiór $M^{<\omega}$. W ZF można dla \emph{ustalonej}
struktury zbiorowej $(M,\in)$ zdefiniować relację spełniania
$\operatorname{Sat}_M(e,\vec p,a)$, gdzie $e$ jest kodem
formuły o wyróżnionej zmiennej $x$, $\vec p$ jest krotką
jej pozostałych parametrów i $a\in M$. Relację tę
konstruujemy przez rekurencję po skończonym drzewie
podformuł: przypadki atomowe używają przynależności i
równości w $M$, spójniki odpowiednich operacji logicznych,
a kwantyfikatory przebiegają zbiór $M$.
Dla każdego kodu $(e,\vec p)$ określamy, przez wyróżnianie,
\[
 A_{e,\vec p}=\{a\in M:\operatorname{Sat}_M(e,\vec p,a)\}.
\]
Schemat zastępowania, stosowany do zbioru dopuszczalnych
kodów i krotek, daje zbiór wszystkich takich wartości
$A_{e,\vec p}$. Jest to dokładnie $\operatorname{Def}(M)$.
Każda wartość jest podzbiorem $M$. Formuły $x\ne x$
i $x=x$ definiują odpowiednio $\varnothing$ i $M$
bez dodatkowych parametrów.
\end{proof}
\begin{definicja}[Hierarchia $L$]\label{def:l-hierarchy}
Przez rekurencję pozaskończoną określamy $L_0=\varnothing$, $L_{\alpha+1}=\operatorname{Def}(L_\alpha)$ oraz $L_\lambda=\bigcup_{\beta<\lambda}L_\beta$ dla ordynałów granicznych $\lambda$. Klasę $L=\bigcup_{\alpha\in\operatorname{Ord}}L_\alpha$ rozumiemy jako klasę wszystkich zbiorów należących do pewnego $L_\alpha$; nie jest to zbiór.
\end{definicja}
\begin{lemat}[Tranzytywność i monotoniczność]\label{lem:l-transitive}
Dla każdego ordynału $\alpha$ zbiór $L_\alpha$ jest tranzytywny i $L_\alpha\subseteq L_{\alpha+1}$.
\end{lemat}
\begin{proof}
Indukcja pozaskończona. Jeśli $L_\alpha$ jest tranzytywny, to dla każdego $a\in L_\alpha$ mamy $a\subseteq L_\alpha$; jest on definiowalny w $L_\alpha$ wzorem $x\in a$ z parametrem $a$, więc $a\in\operatorname{Def}(L_\alpha)$. Każdy element $L_{\alpha+1}$ jest podzbiorem $L_\alpha$, co daje tranzytywność następnego poziomu. Na etapie granicznym $L_\lambda=\bigcup_{\beta<\lambda}L_\beta$ jest tranzytywny. Aby otrzymać $L_\lambda\subseteq L_{\lambda+1}$, zauważmy, że dla $a\in L_\lambda$ mamy $a\subseteq L_\lambda$, a zatem $a$ jest definiowalny nad $L_\lambda$ wzorem $x\in a$ z parametrem $a$.
\end{proof}
\begin{przyklad}[Cztery pierwsze poziomy hierarchii]
\label{ex:l-small-levels-34}
Zaczynamy od $L_0=\varnothing$. Z konwencji dla pustego
uniwersum otrzymujemy $L_1=\{\varnothing\}$.
Struktura o uniwersum $L_1$ ma jeden element, zatem
formuły $x\ne x$ i $x=x$ opisują jej dwa podzbiory:
$\varnothing$ i $L_1$. Stąd
$L_2=\{\varnothing,\{\varnothing\}\}$.
Zbiór $L_2$ ma dwa elementy: zbiór pusty i zbiór
jednoelementowy. W strukturze $(L_2,\in)$ pierwszy
odróżnia formuła $\forall y\ (y\notin x)$,
a drugi jej negacja. Wskutek tego dowolną kombinację
tych dwóch elementów można wskazać formułą, więc
\[
 L_3=\mathcal P(L_2)
 =\{\varnothing,\{\varnothing\},
    \{\{\varnothing\}\},L_2\}.
\]
W szczególności $|L_0|=0$, $|L_1|=1$,
$|L_2|=2$ i $|L_3|=4$.
W istocie dla każdego $n<\omega$ zachodzi $L_n=V_n$.
Uzasadnijmy indukcję. Zbiór $V_n$ jest skończony i tranzytywny.
Każdy jego element $a$ ma jednoznaczny opis bez parametrów:
dla $a=\varnothing$ stosujemy $\varphi_a(x):=\forall y\,(y\notin x)$,
a dla niepustego $a$ (który jest skończony) indukcyjnie
budujemy formułę
\[
 \varphi_a(x):=\forall y\bigl(y\in x\leftrightarrow
        \bigvee_{b\in a}\varphi_b(y)\bigr).
\]
Jest to formuła skończona, ponieważ każdy element $a$
ma skończenie wiele elementów, a relacja przynależności
na zbiorach dziedzicznie skończonych jest dobrze ufundowana.
W strukturze $(V_n,\in)$ formuła $\varphi_a$ rozpoznaje
jednoznacznie $a$ na mocy ekstensjonalności.
Każdy podzbiór $B\subseteq V_n$ jest skończony, więc
definiuje go alternatywa formuł $\varphi_b(x)$ po $b\in B$;
dla pustego $B$ wybieramy $x\ne x$. Stąd
$\operatorname{Def}(V_n)=\mathcal P(V_n)=V_{n+1}$.
Krok $n=0$ daje odrębna konwencja
$\operatorname{Def}(\varnothing)=\{\varnothing\}$.
Ta równość na skończonych poziomach nie oznacza, że
operator $\operatorname{Def}$ ogólnie jest operatorem
zbioru potęgowego; dla nieskończonych struktur
istnieją podzbiory pozbawione definicji z parametrami.
\end{przyklad}

\begin{lemat}[Jawne kodowanie zbiorów dziedzicznie skończonych]
\label{lem:lvi-ackermann-code}
Niech $V_\omega=\bigcup_{n<\omega}V_n$ oznacza zbiór
wszystkich zbiorów dziedzicznie skończonych. Istnieje
bijekcja $a\colon V_\omega\to\omega$ określona
rekurencyjnie wzorem
\[
 a(x)=\sum_{y\in x}2^{a(y)}.
\]
Suma jest skończona, gdyż każdy $x\in V_\omega$
jest zbiorem skończonym. Do konstrukcji i dowodu nie
potrzeba aksjomatu wyboru.
\end{lemat}
\begin{proof}
\emph{Krok 1: istnienie wartości.}
Dla $x\in V_n$ każdy $y\in x$ należy do $V_{n-1}$,
gdy $n>0$. Rekursja po randze pozwala więc obliczyć
wszystkie $a(y)$ przed $a(x)$. Wartość $a(x)$ jest
liczbą naturalną, ponieważ jest sumą skończenie
wielu potęg liczby $2$. Wzór wyznacza wartość jednoznacznie.

\emph{Krok 2: iniekcyjność.}
Przeprowadzamy indukcję po największej z dwóch
skończonych rang zbiorów $x,z$. Jeżeli $a(x)=a(z)$,
to jednoznaczność zapisu liczby naturalnej w systemie
binarnym pokazuje, że
$\{a(y):y\in x\}=\{a(w):w\in z\}$.
Wszystkie elementy $x$ i $z$ mają mniejsze rangi;
z założenia indukcyjnego wynika, że elementy o tym
samym kodzie są równe. Otrzymujemy $x=z$.

\emph{Krok 3: surjekcyjność.}
Stosujemy silną indukcję po $m\in\omega$. Dla $m=0$
mamy $a(\varnothing)=0$. Dla $m>0$ zapis binarny
$m=\sum_{j\in F}2^j$ ma pewien skończony zbiór
pozycji $F\subseteq m$: z $2^j\le m$ wynika $j<m$.
Dla każdego $j\in F$ założenie indukcyjne daje
jedyny $x_j\in V_\omega$ o kodzie $j$.
Zbiór $x=\{x_j:j\in F\}$ należy do $V_\omega$
(skończona rodzina zbiorów dziedzicznie skończonych
ma skończoną sumę rang), a z definicji $a(x)=m$.
To kończy dowód bijektywności.
\end{proof}
\begin{przyklad}[Ręczne obliczenie kodu i pierwszy skok w hierarchii]
\label{ex:lvi-first-infinite-level}
Kodujemy kolejno
$a(\varnothing)=0$, $a(\{\varnothing\})=2^0=1$,
$a(\{\{\varnothing\}\})=2^1=2$ oraz
$a(\{\varnothing,\{\varnothing\}\})=2^0+2^1=3$.
Zatem $L_\omega=V_\omega$ ma jawną numerację
przez liczby naturalne. Każdy element
$L_{\omega+1}=\operatorname{Def}(L_\omega)$
posiada opis składający się z kodu formuły i
skończonej krotki parametrów z $L_\omega$.
Zbiór wszystkich takich opisów jest przeliczalny:
wykorzystujemy kodowanie skończonych krotek
liczb naturalnych. Wobec tego $L_{\omega+1}$
jest przeliczalny. Tymczasem
$\mathcal P(L_\omega)$ nie jest przeliczalny
na mocy twierdzenia Cantora
(Twierdzenie~\ref{thm:zf-cantor}). Otrzymujemy ścisłą
inkluzję
\[
 L_{\omega+1}\subsetneq\mathcal P(L_\omega).
\]
Jest to pierwszy etap, na którym operacja
$\operatorname{Def}$ nie daje całego zbioru
potęgowego. Nie wybieraliśmy żadnych reprezentantów
z nieskończonej rodziny niepustych zbiorów.
\end{przyklad}

\begin{przyklad}[Dlaczego potrzebna jest absolutność]
\label{ex:l-unbounded-quantifier-34}
Ustalmy parametr $x=\varnothing$ i formułę
$\psi(x):=\exists y\ (x\in y)$. W strukturze
$(L_1,\in)$ formuła $\psi(\varnothing)$ jest fałszywa,
gdyż jedynym elementem $L_1$ jest $\varnothing$.
W strukturze $(L_2,\in)$ ta sama formuła jest
prawdziwa: świadkiem jest $y=\{\varnothing\}\in L_2$.
Nieograniczony kwantyfikator egzystencjalny nie jest
więc absolutny nawet między dwoma \emph{tranzytywnymi}
poziomami hierarchii $L$. Natomiast w formule
$\exists y\in a$ elementy $a$ nie zmieniają się
przy przejściu do większego zbioru tranzytywnego.
\end{przyklad}
\section{Absolutność i ostrożność przy kwantyfikatorach}
\begin{definicja}[Formuły ograniczone]\label{def:delta0}
Formułą $\Delta_0$ nazywamy formułę zbudowaną z formuł atomowych przez spójniki i kwantyfikatory ograniczone $\forall x\in y$ i $\exists x\in y$. Formuła jest absolutna między klasami tranzytywnymi, jeśli jej prawdziwość dla parametrów należących do obu klas jest w nich taka sama.
\end{definicja}
\begin{twierdzenie}[Absolutność formuł ograniczonych]\label{thm:delta0-absolute}
Niech $M\subseteq N$ będą tranzytywnymi klasami, w których relacja przynależności jest rzeczywistą relacją $\in$. Dla dowolnej formuły $\Delta_0$ $\varphi$ i dowolnych parametrów $a_1,\ldots,a_n\in M$ zachodzi $M\models\varphi(\vec a)$ wtedy i tylko wtedy, gdy $N\models\varphi(\vec a)$ (przy założeniu, że obie relacje spełniania są rozpatrywane w odpowiedniej metateorii).
\end{twierdzenie}
\begin{proof}
Indukcja po budowie formuły. Równość i przynależność nie zmieniają się przy przejściu z $M$ do $N$. Spójniki zachowują równoważność. Dla $\exists x\in a\,\psi(x,\vec a)$ tranzytywność $M$ zapewnia, że każdy element $a\in M$ również należy do $M$. Zatem kwantyfikator w obu klasach przebiega dokładnie ten sam zbiór elementów $a$. Korzystamy z założenia indukcyjnego dla $\psi$. Analogicznie postępujemy z kwantyfikatorem ogólnym.
\end{proof}
\begin{uwaga}[Granica argumentu]
Zwykły kwantyfikator $\exists x$ może mieć świadków w $N\setminus M$, dlatego powyższe twierdzenie nie dotyczy dowolnych formuł. Pełne twierdzenie $L\models\mathrm{ZFC}+\mathrm{GCH}$
nie wynika jeszcze z samej definicji poziomów hierarchii.
Potrzebne są zasada kondensacji, definiowalny dobry porządek
na klasie $L$ i sprawdzenie schematu zastępowania.
Dowody tych rezultatów znajdują się w dalszych częściach
niniejszego rozdziału (w szczególności w
Lemacie~\ref{lem:l-condensation} i
Twierdzeniu~\ref{thm:l-gch}).
\end{uwaga}

\section{Dlaczego definicja $L$ nie daje jeszcze ZFC?}
\begin{uwaga}[Zbiór a klasa właściwa]
Symbol $L$ oznacza klasę opisaną formułą „istnieje ordynał $\alpha$ taki,
że $x\in L_\alpha$”. Nie wolno stosować do $L$ aksjomatu zbioru
potęgowego tak, jakby był zbiorem. Zdanie $L\models\varphi$
wymaga relatywizacji kwantyfikatorów formuły $\varphi$ do klasy $L$;
nie jest to relacja spełniania dla struktury o zbiorowym uniwersum.
\end{uwaga}
\begin{przyklad}[Definiowalność zależy od struktury]
Niech $M=\{\varnothing\}$. Wtedy $\operatorname{Def}(M)$
składa się z $\varnothing$ i $M$: pierwszy podzbiór
definiuje $x\ne x$, drugi $x=x$. Dla $M=\varnothing$
opisuje go osobna konwencja z
Definicji~\ref{def:constructible-def}, nie relacja spełniania
w strukturze pierwszego rzędu o pustej dziedzinie.
Zgodnie z Przykładem~\ref{ex:l-unbounded-quantifier-34}
różne struktury mogą nadawać różne wartości logiczne
tej samej formule. Dlatego przy operatorze
$\operatorname{Def}$ zawsze wskazujemy uniwersum struktury.
\end{przyklad}

\section{Odbicie formuł do poziomu hierarchii $L$}
\label{sec:l-reflection}
Nie wystarcza powiedzieć, że $L$ jest sumą poziomów $L_\alpha$.
Trzeba jeszcze pokazać, że kwantyfikatory w aksjomatach ZF
można kontrolować wewnątrz tej klasy. Poniższy lemat
jest narzędziem do tego celu.

\begin{lemat}[Odbicie skończonej rodziny formuł]
\label{lem:l-finite-reflection}
Niech $\Phi$ będzie skończonym zbiorem formuł pierwszego
rzędu języka $\{\in\}$, domkniętym na podformuły.
Dla każdego ordynału $\eta$ istnieje ordynał graniczny
$\theta>\eta$ taki, że dla każdej $\varphi\in\Phi$
i wszystkich parametrów $\vec a\in L_\theta$ zachodzi
\[
(L_\theta,\in)\models\varphi(\vec a)
\quad\Longleftrightarrow\quad L\models\varphi(\vec a).
\]
Przez $L\models\varphi$ rozumiemy formułę otrzymaną
przez ograniczenie wszystkich kwantyfikatorów do
\emph{definiowalnej klasy} $L$; nie zakładamy istnienia
zbiorowej struktury o uniwersum $L$.
\end{lemat}
\begin{proof}
\emph{Etap 1: wybieranie etapów, nie świadków.}
Ustalmy formułę $\exists y\,\psi(y,\vec x)$ będącą
podformułą którejś formuły z $\Phi$. Jeżeli
$\vec a\in L_\alpha$ i $L\models\exists y\,\psi(y,\vec a)$,
to istnieje najmniejszy ordynał $\beta$ taki, że
$\exists y\in L_\beta\,\psi^L(y,\vec a)$.
Ta liczba porządkowa jest wyznaczona jednoznacznie
przez $\vec a$ i kod formuły. Ponieważ krotek
parametrów z $L_\alpha$ jest tylko zbiór,
aksjomat zastępowania w metateorii ZF zapewnia
wspólne ograniczenie dla ich najmniejszych etapów
oraz dla skończenie wielu formuł egzystencjalnych.
Nie wybieramy tutaj naraz elementów $y$ z całej rodziny:
wystarcza zebrać \emph{ordynały} ich pierwszego wystąpienia.

\emph{Etap 2: jednoznacznie określone domknięcie.}
Aby nie używać wyboru przeliczalnego,
ustalamy $\alpha_0:=\eta+1$.
Dla każdej $\alpha$ niech $E_\alpha$
będzie zbiorem par $(e,\vec a)$,
gdzie $e$ jest kodem podformuły
egzystencjalnej z $\Phi$,
$\vec a$ jest odpowiednią krotką
parametrów z $L_\alpha$ i w $L$
istnieje jej świadek.
Zbiór $E_\alpha$ istnieje przez
separację w skończonej sumie
zbiorów krotek parametrów.
Każdej parze z $E_\alpha$
przyporządkowujemy najmniejszy
ordynał $\beta_{e,\vec a}$ taki,
że świadek należy do
$L_{\beta_{e,\vec a}}$.
Zbiór wszystkich tych ordynałów
istnieje przez zastępowanie.
Kładziemy jednoznacznie
\[
F(\alpha):=
\sup\bigl(\{\alpha+1\}
 \cup\{\beta_{e,\vec a}+1:
                  (e,\vec a)\in E_\alpha\}\bigr).
\]
Definiujemy przez rekurencję na $\omega$
$\alpha_{n+1}:=F(\alpha_n)$.
Mamy $\alpha_{n+1}>\alpha_n$
i każdy wskazany świadek należy do
$L_{\alpha_{n+1}}$.
Niech $\theta:=\sup_{n\in\omega}\alpha_n$;
jest to ordynał graniczny większy od $\eta$.
Dla każdej skończonej krotki
$\vec a\in L_\theta^{<\omega}$
istnieje $n$ takie, że
wszystkie jej współrzędne należą
do $L_{\alpha_n}$.
Jeżeli odpowiednia podformuła
egzystencjalna ma świadka w $L$,
to ma go już w
$L_{\alpha_{n+1}}\subseteq L_\theta$.
Wszystkie etapy rekurencji
wyznaczono jednoznacznymi formułami ZF;
nie zastosowano AC ani DC.

\emph{Etap 3: indukcja po formułach.}
Dla formuł atomowych prawdziwość jest taka sama,
bo $L_\theta$ jest tranzytywny. Spójniki zachowują
równoważność. Gdy $\exists y\,\psi(y,\vec a)$ jest
prawdziwe w $L$, Etap 2 dostarcza świadka
$y\in L_\theta$, a założenie indukcyjne daje
prawdziwość $\psi$ w $L_\theta$.
W przeciwną stronę świadek z $L_\theta$ jest
również elementem $L$, więc stosujemy
założenie indukcyjne. Kwantyfikator ogólny
wyrażamy przez negację egzystencjalnego.
Otrzymujemy jednocześnie równoważność dla wszystkich
formuł z $\Phi$.
\end{proof}

\section{Dowód, że $L$ spełnia aksjomaty ZF}
\label{sec:l-model-zf}
\begin{twierdzenie}[Uniwersum konstruowalne jest modelem wewnętrznym ZF]
\label{thm:l-model-zf}
W metateorii ZF klasa $L$ jest tranzytywna, zawiera
wszystkie ordynały oraz spełnia każdy aksjomat ZF
po relatywizacji kwantyfikatorów do $L$.
\end{twierdzenie}
\begin{proof}
\emph{Krok 1: tranzytywność, ordynały i ekstensjonalność.}
Tranzytywność wynika z Lematu~\ref{lem:l-transitive}.
Uzasadnijmy dokładniej równość
$L_\alpha\cap\operatorname{Ord}=\alpha$,
która jest istotna dla wszystkich dalszych etapów.
W ZF zbiór $x$ jest ordynałem wtedy i tylko wtedy,
gdy $x$ jest przechodni i każdy element $x$
jest przechodni; dzięki ufundowaniu
jest to warunek równoważny dobremu
uporządkowaniu przez $\in$.
Warunek przechodniości można zapisać
\[
\forall u\in x\;\forall v\in u\;(v\in x),
\]
a warunek przechodniości każdego elementu
$x$ ma analogiczną postać
z kwantyfikatorami ograniczonymi.
Zatem predykat $\operatorname{Ord}(x)$
ma definicję $\Delta_0$ i jest
absolutny dla zbiorów przechodnich
(Twierdzenie~\ref{thm:delta0-absolute}).
Przez indukcję po $\alpha$, gdy
$L_\alpha\cap\operatorname{Ord}=\alpha$,
zbiór
\[
\alpha=\{x\in L_\alpha:
             (L_\alpha,\in)\models\operatorname{Ord}(x)\}
\]
jest definiowalnym podzbiorem $L_\alpha$.
Stąd $\alpha\in L_{\alpha+1}$,
a ponieważ każdy element $L_{\alpha+1}$
jest podzbiorem przechodniego $L_\alpha$,
nie może pojawić się na tym etapie
ordynał większy od $\alpha$:
jeśli $\gamma\in L_{\alpha+1}$
jest ordynałem, to
$\gamma\subseteq L_\alpha$ i
$\gamma\subseteq\alpha$, zatem
$\gamma\le\alpha$.
Na etapach granicznych korzystamy
z monotoniczności i sumy poziomów,
a dla $\alpha=0$ obie strony są puste.
To kończy indukcję i pokazuje,
że wszystkie ordynały należą do $L$.
Równość i przynależność są rzeczywistymi relacjami
na zbiorach, więc ekstensjonalność zachodzi w $L$.

\emph{Krok 2: pustka, pary, suma i nieskończoność.}
Dla $a,b\in L$ wybieramy $\alpha$ z $a,b\in L_\alpha$.
Zbiór $\{a,b\}$ jest podzbiorem $L_\alpha$
definiowalnym w $L_\alpha$ formułą $x=a\lor x=b$;
stąd $\{a,b\}\in L_{\alpha+1}$. Ponieważ
$L_\alpha$ jest tranzytywny, suma $\bigcup a$
jest podzbiorem $L_\alpha$ definiowalnym przez
$\exists y\in a\,x\in y$. Zatem $\bigcup a\in L$.
Zbiór pusty należy do $L_1$.
Wykażmy indukcyjnie $L_n=V_n$ dla $n\in\omega$.
Dla $n=0$ oba zbiory są puste.
Jeśli $L_n=V_n$ jest skończony,
to każdy jego podzbiór $B$ jest
definiowalny nad $L_n$: dla
$B=\{a_1,\ldots,a_k\}$ wystarczy
formuła
$x=a_1\lor\cdots\lor x=a_k$
z parametrami $a_i\in L_n$,
a dla $B=\varnothing$ formuła $x\ne x$.
Zatem
$\operatorname{Def}(L_n)=\mathcal P(L_n)$
(również dla $n=0$ na mocy konwencji),
czyli $L_{n+1}=V_{n+1}$.
Stąd $L_\omega=V_\omega$ jest
zbiorem wszystkich dziedzicznie
skończonych zbiorów.
Z wcześniej uzasadnionej
absolutności predykatu
$\operatorname{Ord}$ wynika,
że zbiór jego elementów będących
skończonymi ordynałami jest
definiowalny nad $L_\omega$
i równy $\omega$.
Stąd $\omega\in L_{\omega+1}$,
co daje aksjomat nieskończoności.

\emph{Krok 3: wyróżnianie.}
Niech $a,\vec p\in L$ i niech $\varphi(x,\vec p)$
będzie ustaloną formułą. Wybieramy $\theta$ z
Lematu~\ref{lem:l-finite-reflection} dla tej formuły
oraz wszystkich jej podformuł, tak duży, aby
$a,\vec p\in L_\theta$. Wtedy
\[
\{x\in a:L\models\varphi(x,\vec p)\}
=\{x\in L_\theta:x\in a\land
L_\theta\models\varphi(x,\vec p)\}.
\]
Prawa strona jest definiowalnym podzbiorem $L_\theta$,
a więc należy do $L_{\theta+1}$.

\emph{Krok 4: zbiór potęgowy.}
Ustalmy $a\in L$. Zewnętrzna rodzina
$B=\{b\subseteq a:b\in L\}$ jest zbiorem w ZF:
wydzielamy ją z $\mathcal P(a)$ za pomocą formuły
„$b\in L$”. Dla każdego $b\in B$ istnieje najmniejszy
etap $\beta_b$ z $b\in L_{\beta_b}$; przez
zastępowanie istnieje wspólne ograniczenie tych
etapów. Wybieramy $\theta$ tak duże, że
$B\subseteq L_\theta$ oraz $a\in L_\theta$.
Wtedy
\[
B=\{b\in L_\theta:\forall x\in b\ (x\in a)\}.
\]
Zbiór $B$ jest więc definiowalny nad $L_\theta$,
a zatem $B\in L_{\theta+1}$. Wewnątrz $L$
pełni dokładnie rolę zbioru potęgowego $a$:
zawiera wszystkie i tylko podzbiory $a$, które
należą do $L$. Nie należy go utożsamiać z całym
zewnętrznym $\mathcal P(a)$.

\emph{Krok 5: zastępowanie.}
Załóżmy, że $L\models\forall x\in a\ \exists!y\,
\varphi(x,y,\vec p)$, gdzie $a,\vec p\in L$.
W metateorii ZF formuła $\varphi^L$ jest zwykłą
formułą pierwszego rzędu z klasą $L$ określoną
definiowalnie. Dlatego przez zewnętrzny schemat
zastępowania istnieje zbiór $B$ wszystkich jej
jednoznacznych wartości $y$ dla $x\in a$.
Każda z tych wartości jest w $L$. Wybieramy
$\theta$ tak duże, aby $a,B,\vec p\subseteq L_\theta$
(w przypadku $a$ i parametrów chodzi także o ich
przynależność do $L_\theta$), a ponadto aby
$L_\theta$ poprawnie rozstrzygał $\varphi$
na swoich parametrach, zgodnie z Lematem
\ref{lem:l-finite-reflection}. Wówczas
\[
B=\{y\in L_\theta:\exists x\in a\;
L_\theta\models\varphi(x,y,\vec p)\}.
\]
Dzięki odbiciu nie pojawiają się tu fałszywe wartości.
Prawa strona jest definiowalna nad $L_\theta$;
wobec tego $B\in L_{\theta+1}$. Dowiedliśmy
zastępowania \emph{wewnątrz $L$}.

\emph{Krok 6: ufundowanie.}
Jeśli $a\in L$ jest niepusty w $L$, to jest niepusty
w zewnętrznym ZF. Aksjomat ufundowania daje element
$\in$-minimalny $x\in a$. Tranzytywność $L$ oznacza,
że $x\in L$, a więc jest on również świadkiem
ufundowania w $L$. Wszystkie aksjomaty ZF zostały
sprawdzone; nie wykorzystaliśmy AC w metateorii.
\end{proof}


\begin{lemat}[Absolutność hierarchii konstruowalnej wewnątrz $L$]
\label{lem:l-hierarchy-internal}
Niech $\alpha$ będzie dowolną liczbą porządkową. Poziom hierarchii
konstruowalnej obliczony wewnątrz klasy $L$ jest tym samym zbiorem,
co poziom $L_\alpha$ obliczony w otaczającej teorii ZF:
\[
 (L_\alpha)^L=L_\alpha.
\]
W szczególności, jeżeli $M=L_\alpha$, to dla każdej \emph{ustalonej}
formuły pierwszego rzędu $\varphi$ i parametrów z $M$ relacja
spełniania $M\models\varphi$ jest jednakowo obliczana w $L$ i w ZF.
\end{lemat}
\begin{proof}
Najpierw ustalmy $M\in L$. Prawdziwość formuły atomowej
w zbiorowej strukturze $(M,\in)$ zależy tylko od równości
i przynależności jej elementów. Dla spójników prawdziwość
wyznacza się z wartości logicznych podformuł; dla kwantyfikatora
$\exists x$ przebiega się dokładnie zbiór $M$, który jest ten sam
w klasie tranzytywnej $L$ i w otaczającej teorii. Indukcja
po budowie \emph{ustalonej} formuły daje więc absolutność
relacji spełniania w $(M,\in)$.

Ponadto, jeżeli $B\subseteq M$ jest podzbiorem definiowalnym
w $(M,\in)$ za pomocą formuły i parametrów z $M$, to
$B\in L$: można go wydzielić w $L$ ze zbioru $M$ formułą
wyrażającą relację spełniania zbiorowej struktury.
Odwrotnie, każda definicja obliczona w $L$ opisuje dokładnie
ten sam podzbiór $M$ co w ZF. Zatem
$\operatorname{Def}^L(M)=\operatorname{Def}(M)$.

Przeprowadzamy teraz indukcję pozaskończoną po $\alpha$.
Dla zera oba poziomy są puste. Jeżeli mają tę samą wartość
na etapie $\alpha$, to w kroku następnym wyznacza je ten sam
operator $\operatorname{Def}$. Na etapie granicznym obie
hierarchie biorą sumę tej samej rodziny wcześniejszych
poziomów. Każda liczba porządkowa należy do $L$, więc
indukcja obejmuje wszystkie liczby porządkowe zewnętrznego
uniwersum.
\end{proof}

\section{Kanoniczny dobry porządek $L$ i aksjomat wyboru}
\label{sec:l-choice}
\begin{twierdzenie}[Definiowalny dobry porządek uniwersum konstruowalnego]
\label{thm:l-definable-wellorder}
Istnieje formuła bez parametrów definiująca w $L$
relację $<_{L}$, której ograniczenie do każdego
zbioru należącego do $L$ jest dobrym porządkiem.
W szczególności $L\models\mathrm{AC}$.
\end{twierdzenie}
\begin{proof}
\emph{Etap 1: porządkujemy zbiór definicji.}
Jeśli $L_\alpha$ jest dobrze uporządkowany,
każdy element $L_{\alpha+1}=\operatorname{Def}(L_\alpha)$
ma co najmniej jeden opis: numer formuły pierwszego
rzędu i skończoną krotkę parametrów z $L_\alpha$.
Porządkujemy kody, porównując najpierw długość
krotki i numer formuły, a następnie leksykograficznie
parametry. Dla ustalonej długości iloczyn skończonej
liczby dobrych porządków z porządkiem leksykograficznym
jest dobrym porządkiem; suma po długościach i numerach
też jest dobrze uporządkowana. Wybieramy dla każdego
zbioru najwcześniejszy spośród jego kodów.

\emph{Etap 2: porządkujemy całe poziomy spójnie.}
Przez rekurencję pozaskończoną tworzymy porządek
$<_{L,\alpha}$ na $L_\alpha$. Na etapie następnym
zachowujemy dotychczasowy porządek na $L_\alpha$,
a \emph{nowe} elementy $L_{\alpha+1}\setminus L_\alpha$
umieszczamy za starymi zgodnie z Etapem 1.
Na etapie granicznym łączymy porządki z etapów
wcześniejszych. Każdy niepusty zbiór konstruktybilny
ma element pojawiający się na najwcześniejszym etapie;
w tej warstwie istnieje najmniejszy element według
wybranego porządku. Otrzymujemy $<_{L}$.
Cała rekurencja jest jednolicie definiowalna z
formuły kodującej $\alpha\mapsto L_\alpha$,
bez pozaskończonego ciągu arbitralnych wyborów.

\emph{Etap 3: wyprowadzamy AC.}
Dla każdego $a\in L$ relacja $<_{L}$ ograniczona do
$a$ jest dobrym porządkiem, a z aksjomatów ZF
spełnionych w $L$ jego wykres jest zbiorem w $L$
(otrzymujemy go przez wyróżnianie w $a\times a$).
Z twierdzenia o dobrym uporządkowaniu Zermela,
Twierdzenie~\ref{thm:choice-equivalence},
wynika aksjomat wyboru w $L$.
\end{proof}

\section{Lemat kondensacji i hipoteza continuum uogólniona}
\label{sec:l-condensation-gch}
\begin{definicja}[Kolaps Mostowskiego]
\label{def:l-mostowski-collapse}
Dla ekstensjonalnej i dobrze ufundowanej relacji
$\in$ na zbiorze $X$ kolapsem Mostowskiego nazywamy
funkcję $\pi:X\to\bar X$ zadaną rekurencją
\[
\pi(x)=\{\pi(y):y\in X\land y\in x\}.
\]
Jej obraz $\bar X$ jest tranzytywny, a $\pi$ jest
izomorfizmem $(X,\in)$ na $(\bar X,\in)$.
Gdy $X$ jest podstrukturą elementarną pewnego
zbioru tranzytywnego, relacja jest ekstensjonalna
na $X$ dzięki elementarności.
\end{definicja}

\begin{lemat}[Elementarność przekroju z elementem podstruktury]
\label{lem:l-elementary-section-75}
Niech $D$ będzie niepustym zbiorem tranzytywnym,
$X\preccurlyeq(D,\in)$,
a $A\in X$ będzie niepustym elementem $D$.
Wówczas
\[
(X\cap A,\in)\preccurlyeq(A,\in).
\]
Jeżeli $\pi:X\to N$ jest kolapsem Mostowskiego, to
$\pi\restriction(X\cap A)$ jest izomorfizmem
$(X\cap A,\in)$ na $(\pi(A),\in)$.
\end{lemat}
\begin{proof}
Ustalmy formułę $\varphi(u,\vec v)$ języka teorii mnogości
i parametry $\vec a\in(X\cap A)^n$.
Prawdziwość $\varphi$ w strukturze $(A,\in)$ jest
równoważna prawdziwości relatywizacji $\varphi^A$
w $(D,\in)$, ponieważ $A\in D$ i tranzytywność $D$ zapewnia,
że wszystkie elementy $A$ należą do uniwersum $D$.
Jeżeli $(A,\in)\models\exists u\,\varphi(u,\vec a)$,
to $(D,\in)$ spełnia formułę
$\exists u\in A\,\varphi^A(u,\vec a)$.
Wszystkie parametry tej formuły należą do $X$.
Elementarność $X\preccurlyeq D$ daje świadka
$u\in X$ takiego, że $u\in A$ i
$(A,\in)\models\varphi(u,\vec a)$.
Kryterium Tarskiego--Vaughta dowodzi pierwszej tezy.
W drugiej części definicja kolapsu daje
$\pi(A)=\{\pi(u):u\in X\cap A\}$, a zachowanie
relacji $\in$ i różnowartościowość $\pi$ zapewniają
wskazany izomorfizm.
\end{proof}

\begin{uwaga}[Co wynika, a co nie wynika z elementarności]
\label{rem:l-sections-corrected-75}
Z założeń $X\preccurlyeq L_\theta$ i $L_\gamma\in X$
\emph{rzeczywiście wynika}
$(X\cap L_\gamma,\in)\preccurlyeq(L_\gamma,\in)$,
bez dodatkowej hipotezy. Nie wynika natomiast od razu
$\pi(L_\gamma)=L_{\pi(\gamma)}$:
prawa strona wymaga przeprowadzenia indukcji po poziomach.
Tych dwóch stwierdzeń nie należy utożsamiać.
\end{uwaga}

\begin{przyklad}[Dlaczego tranzytywność struktury większej jest potrzebna]
\label{ex:l-elementary-section-transitivity-75}
Niech $A=\{\varnothing\}$, $D=\{A\}$ i $X=D$.
Wtedy $X\preccurlyeq(D,\in)$ jest tautologią
(struktura jest identyczna z własną podstrukturą),
$A\in X$ i $A$ jest niepusty. Jednak
$X\cap A=\varnothing$, ponieważ
$\varnothing\notin D$. Nie otrzymujemy zatem
podstruktury elementarnej niepustej struktury
$(A,\in)$. Przyczyną jest nietranzytywność $D$;
Lemat~\ref{lem:l-elementary-section-75}
celowo zakłada tranzytywność większego uniwersum.
\end{przyklad}

\begin{lemat}[Kolaps Mostowskiego a zbiory definiowalne]
\label{lem:l-collapse-def-compatibility}
Niech $\theta>0$ będzie graniczną liczbą porządkową,
$X\preccurlyeq(L_\theta,\in)$,
$\pi:X\to N$ kolapsem Mostowskiego, a $\gamma\in X$
liczbą porządkową taką, że $L_\gamma\in X$.
Jeżeli dodatkowo
$\pi(L_\gamma)=L_{\pi(\gamma)}$, to
\[
\pi[X\cap\operatorname{Def}(L_\gamma)]
  =\operatorname{Def}(L_{\pi(\gamma)}).
\]
Nie zakłada się osobno elementarności przekroju
$X\cap L_\gamma$: wynika ona z
Lematu~\ref{lem:l-elementary-section-75} dla $\gamma>0$.
\end{lemat}
\begin{proof}
Jeżeli $\gamma=0$, to
$L_0=\varnothing$, $\operatorname{Def}(L_0)=\{\varnothing\}$,
a $0\in X$ i $\pi(0)=0$; teza zachodzi.
Załóżmy więc $\gamma>0$ i oznaczmy
$A=L_\gamma$, $A'=L_{\pi(\gamma)}$.
Lemat~\ref{lem:l-elementary-section-75} daje
$(X\cap A,\in)\preccurlyeq(A,\in)$ oraz izomorfizm
$\pi\restriction(X\cap A)$ na $(A',\in)$.
Dla każdej \emph{ustalonej} formuły $\varphi(u,\vec v)$
jej prawdziwość z parametrami w $X\cap A$ jest więc
taka sama w $A$ i w $A'$ po przeniesieniu parametrów
przez $\pi$. Nie stosujemy predykatu prawdy dla klasy $L$.

\emph{Pierwsza inkluzja.}
Niech $z\in X\cap\operatorname{Def}(A)$.
W metateorii istnieje formuła $\varphi$ i krotka
$\vec p\in A^{<\omega}$ takie, że
$z=\{u\in A:(A,\in)\models\varphi(u,\vec p)\}$.
Po \emph{ustaleniu} tej formuły zdanie określające $z$
jest pierwszego rzędu w $L_\theta$:
kwantyfikatory formuły $\varphi$ relatywizujemy do $A$.
Ponieważ $A,z\in X$, elementarność zapewnia
krotkę $\vec p'\in(X\cap A)^{<\omega}$ definiującą
\emph{ten sam} zbiór $z$.
Dla $v\in A'$ wybierzmy $u\in X\cap A$ z
$\pi(u)=v$. Wówczas
\begin{align*}
 v\in\pi(z)
 &\iff u\in z
 \iff (A,\in)\models\varphi(u,\vec p')\\
 &\iff (A',\in)\models
            \varphi(v,\pi(\vec p')).
\end{align*}
Stąd $\pi(z)\in\operatorname{Def}(A')$.

\emph{Druga inkluzja.}
Niech $b=\{v\in A':(A',\in)\models
\varphi(v,\vec q)\}$, gdzie $\vec q\in(A')^{<\omega}$.
Wybieramy preobrazy $\vec p\in(X\cap A)^{<\omega}$
parametrów $\vec q$ i definiujemy
$z=\{u\in A:(A,\in)\models\varphi(u,\vec p)\}$.
Ponieważ $z\in\operatorname{Def}(A)=L_{\gamma+1}$,
a $\gamma+1<\theta$, mamy $z\in L_\theta$.
Zbiór $z$ jest \emph{jednoznacznie} określony w
$L_\theta$ przez $A,\vec p$ i ustaloną relatywizację
$\varphi^A$; elementarność daje $z\in X$.
Poprzednie równoważności pokazują $\pi(z)=b$.
To kończy dowód obu inkluzji.
\end{proof}


\begin{lemat}[Tablica spełniania pojedynczej formuły i jej absolutność]
\label{lem:l-satisfaction-table-transitive-73}
Niech $A$ będzie niepustym zbiorem i niech $e\in\omega$
będzie kodem ustalonej formuły pierwszego rzędu $\varphi_e$.
Istnieje jednoznacznie wyznaczony zbiór $T_{A,e}$,
który koduje, dla każdej podformuły $\psi$ formuły
$\varphi_e$ i każdego wartościowania jej zmiennych
w $A$, prawdziwość $\psi$ w $(A,\in)$.
Jeżeli $N$ jest tranzytywnym \emph{modelem ZF}
zawierającym $A$ i $e$, to tablica skonstruowana w $N$
jest dokładnie $T_{A,e}$ obliczoną w otaczającym uniwersum.
W szczególności prawdziwość każdej ustalonej formuły
w zbiorowej strukturze $(A,\in)$ jest absolutna między
takimi modelami a otaczającym uniwersum.
\end{lemat}
\begin{proof}
Niech $W$ będzie skończonym zbiorem zmiennych występujących
w $\varphi_e$, również tych związanych, a $\operatorname{Sub}(e)$
-- skończonym zbiorem kodów podformuł. Zbiór
$A^W$ wszystkich funkcji $s:W\to A$ istnieje w ZF.
Poszukujemy podzbioru
$T\subseteq\operatorname{Sub}(e)\times A^W$,
którego elementy $(\psi,s)$ oznaczają, że
$(A,\in)\models\psi[s]$.
Dla formuły atomowej $v\in w$ lub $v=w$ prawdziwość
zależy wyłącznie od tego, czy odpowiednie wartości
$s(v),s(w)\in A$ pozostają w relacji przynależności
lub są równe. Dla negacji bierzemy dopełnienie
odpowiedniego przekroju $A^W$; dla koniunkcji
-- przecięcie dwóch takich przekrojów. Jeżeli
$\psi$ ma postać $\exists v\,\chi$, przyjmujemy
$(\psi,s)\in T$ wtedy i tylko wtedy, gdy istnieje
$a\in A$ takie, że $(\chi,s[v\mapsto a])\in T$.
Jest to kwantyfikator po \emph{ustalonym zbiorze} $A$.
Indukcja po wysokości drzewa składniowego konstruuje
wszystkie przekroje $T$ za pomocą wyróżniania,
operacji zbiorowych i skończenie wielu etapów.
Ta sama indukcja daje jednoznaczność: wartości dla
formuł atomowych są ustalone, a każdy konstruktor
logiczny jednoznacznie określa wartość na podstawie
ściśle mniejszych podformuł.

Jeżeli $A,e\in N$ i $N$ jest tranzytywnym modelem ZF,
to $W$, $\operatorname{Sub}(e)$ i $A^W$ są w $N$
i mają te same elementy co w uniwersum zewnętrznym:
funkcje o skończonej dziedzinie można zakodować przez
skończone zbiory par, a tranzytywność zapewnia
zgodność ich składników. Wszystkie wymienione
operacje logiczne i kwantyfikacja po $A$ są absolutne.
Indukcja po tym samym skończonym drzewie pokazuje,
że obie konstrukcje przypisują każdemu $(\psi,s)$
tę samą wartość. Wobec jednoznaczności tablice są równe.
\end{proof}
\begin{uwaga}[Granica zastosowania lematu o tablicy]
\label{rem:l-table-vs-stage-73}
Lemat~\ref{lem:l-satisfaction-table-transitive-73}
\emph{nie} pozwala zastąpić w jego założeniach
tranzytywnego modelu ZF dowolnym poziomem $L_\theta$
ani dowolną podstrukturą elementarną $(L_\theta,\in)$.
Taki poziom nie musi spełniać całej ZF, a istnienie
\emph{jednocześnie wszystkich} tablic dla nieograniczonej
rodziny kodów formuł nie wynika ze skończoności
\emph{jednej} tablicy. Dlatego w dowodzie ogólnego
lematu kondensacji nadal potrzebny jest odrębny
argument o jednolitych historiach poziomów $L$.
\end{uwaga}

\begin{uwaga}[Dlaczego kod formuły nie jest predykatem prawdy dla $L$]
\label{rem:l-finite-satisfaction-in-limit}
W Lemacie~\ref{lem:l-collapse-def-compatibility} symbol
$\operatorname{Sat}(L_\gamma,e,\vec p,z)$ opisuje
prawdziwość \emph{jednej formuły o kodzie $e$ w
zbiorowej strukturze $L_\gamma$}. Nie jest to
jednolity predykat prawdy dla klasy właściwej $L$.
Można formalnie obliczać prawdziwość tej jednej formuły
przez rekurencję po jej \emph{skończonym} drzewie
podformuł: w kroku atomowym używamy relacji $\in$,
a każdy kolejny krok polega na operacji logicznej
albo na kwantyfikowaniu po zbiorze $L_\gamma$.
Jeżeli kod $e$ jest ustalony, tablice wyników tej
skończonej rekurencji można zakodować w pewnym poziomie
$L_{\gamma+n(e)}$, gdzie $n(e)<\omega$ zależy od
budowy formuły. Dla granicznego $\theta>\gamma$
zachodzi $\gamma+n(e)<\theta$.
Tę \emph{skończoną poprawność} wykorzystuje argument
o kolapsie; nie zakładamy, że cała nieskończona
relacja spełniania dowolnych formuł jest elementem
$L_\theta$.
\end{uwaga}



\begin{lemat}[Lokalna tablica spełniania w granicznym poziomie $L$]
\label{lem:l-local-satisfaction-bound-76}
Niech $\varphi$ będzie \emph{ustaloną} formułą pierwszego
rzędu języka $\{\in\}$, a $W$ skończonym zbiorem
wszystkich jej zmiennych. Istnieje liczba naturalna
$n_\varphi>0$ zależna wyłącznie od składni $\varphi$
taka, że dla każdej liczby porządkowej $\gamma>0$
zbiór
\[
T_{\gamma,\varphi}=\bigl\{(\psi,s):
  \psi\text{ jest podformułą }\varphi,\ 
  s:W\to L_\gamma,\ (L_\gamma,\in)\models\psi[s]\bigr\}
\]
należy do $L_{\gamma+n_\varphi}$.
W szczególności, jeżeli $\theta>0$ jest graniczny
i $0<\gamma<\theta$, to $T_{\gamma,\varphi}\in L_\theta$.
Nie twierdzimy, że istnieje jedna liczba $n$
wspólna dla \emph{wszystkich} formuł.
\end{lemat}
\begin{proof}
Niech $A=L_\gamma$. Ponieważ $\gamma>0$, zbiór
$A$ jest niepusty i $A\in L_{\gamma+1}$.
Zbiór funkcji $A^W$ otrzymujemy ze skończonej
liczby iloczynów kartezjańskich i ograniczeń do
krotek reprezentujących funkcje o dziedzinie $W$.
Liczba potrzebnych operacji zależy od $|W|$, nie
od liczby elementów $A$. Kody zmiennych, podformuł
i ich drzewa składniowego są dziedzicznie skończone.

Budujemy osobno dla każdej podformuły $\psi$ relację
$S_\psi=\{s\in A^W:(A,\in)\models\psi[s]\}$.
Dla równości i przynależności relacja taka powstaje
przez przecięcie zbioru wartościowań z relacją
atomową. Dla negacji używamy różnicy
$A^W\setminus S_\chi$, a dla koniunkcji — przecięcia.
Jeśli $\psi=\exists v\,\chi$, to
\[
S_\psi=\{s\in A^W:\exists a\in A\ 
                   s[v\mapsto a]\in S_\chi\}.
\]
Zbiór par $(s,a)$ z taką własnością otrzymujemy
przez skończone operacje na produktach $A^W\times A$;
następnie rzutujemy na współrzędną $s$. Jest to
skończone złożenie operacji Gödla z
Lematu~\ref{lem:l-godel-normal-closure}.
W każdym kroku korzystamy wyłącznie z wyników
odpowiadających \emph{ściśle mniejszym} podformułom.

Każda z użytych operacji ma stałą realizację za pomocą
formuł ograniczonych, iloczynów, różnic i rzutów.
Jeżeli jej argumenty należą do $L_\delta$, to wynik
należy do pewnego $L_{\delta+c}$, gdzie stała skończona
$c$ zależy od \emph{rodzaju} operacji, a nie od
$\delta$ ani liczności jej argumentów. Istotnie,
$L_\delta$ jest tranzytywny, zaś odpowiednie pośrednie
zbiory otrzymuje się z $L_\delta$ przez ustaloną liczbę
operacji $\operatorname{Def}$; równoważnie można
wykorzystać normalną postać z przywołanego lematu.
W drzewie $\varphi$ jest tylko skończenie wiele węzłów.
Sumując liczbę operacji i koszt kodowania par
$(\psi,s)$, wybieramy jedno $n_\varphi$ tak, żeby
wszystkie potrzebne relacje i ich skończona suma
$T_{\gamma,\varphi}$ należały do $L_{\gamma+n_\varphi}$.

Dla granicznego $\theta$ i $\gamma<\theta$ każdy
skończony następnik $\gamma+n_\varphi$ jest
mniejszy od $\theta$, więc tablica należy do
$L_\theta$. Argument nie stosuje zastępowania
wewnątrz $L_\theta$ i nie zbiera w $L_\theta$
tablic wszystkich formuł naraz.
\end{proof}

\begin{przyklad}[Dwa kwantyfikatory i granica jednolitości]
\label{ex:l-local-satisfaction-76}
Niech $\varphi(u)$ będzie formułą
$\exists v\,(u\in v\land\forall w\,(w\in v\Rightarrow w=u))$.
Dla $A=L_\gamma$ i $\gamma>0$ budujemy najpierw
relację atomową $u\in v$ na $A^2$, następnie
relację $w\in v\Rightarrow w=u$ na $A^3$.
Jej dopełnienie i rzut względem $w$ opisują
kwantyfikator $\forall w$; przecięcie z relacją
$u\in v$ i rzut względem $v$ opisują
kwantyfikator $\exists v$. Każdy krok jest skończoną
operacją na zbiorach wartościowań. Poziom konstruktywny
wystarczający do uzyskania całej tablicy zależy
od składni tej jednej formuły, a nie od $\gamma$.
Zmiana $\varphi$ na formuły o rosnącej liczbie
zagnieżdżonych kwantyfikatorów zmienia ten koszt;
nie wolno więc wyciągać wniosku o jednej wspólnej
tablicy wszystkich formuł w dowolnym $L_\theta$.
\end{przyklad}


\begin{lemat}[Operacje Gödla i postać normalna definicji]
\label{lem:l-godel-normal-closure}
Dla zbiorów $A,B$ definiujemy następujące operacje
($G_6,G_7,G_8,G_9,G_{10}$ są jednoargumentowe):
\begin{align*}
G_1(A,B)&=\{A,B\}, &G_2(A,B)&=A\times B,\\
G_3(A,B)&=\{(u,v)\in A\times B:u\in v\},
 &G_4(A,B)&=A\setminus B,\\
G_5(A,B)&=A\cap B, &G_6(A)&=\bigcup A,\\
G_7(A)&=\operatorname{dom}(A),
 &G_8(A)&=\{(v,u):(u,v)\in A\},\\
G_9(A)&=\{(u,w,v):(u,v,w)\in A\},
 &G_{10}(A)&=\{(w,u,v):(u,v,w)\in A\}.
\end{align*}
Krotki w ostatnich dwóch wzorach są kodowane jednolicie.
Każda z relacji $B=G_i(\vec A)$ jest wyrażalna
formułą $\Delta_0$; w szczególności jej wartość jest
absolutna dla zbiorów tranzytywnych, w których istnieje.
Jeżeli $A$ jest tranzytywny, to
\[
\operatorname{Def}(A)
=\operatorname{cl}_{G_1,\ldots,G_{10}}(A\cup\{A\})
\cap\mathcal P(A),
\]
gdzie $\operatorname{cl}_{G_1,\ldots,G_{10}}$ oznacza
zbiór wszystkich wyników \emph{skończonych} złożeń tych
operacji z argumentami z podanego zbioru.
\end{lemat}
\begin{proof}
\emph{Absolutność.} Na przykład relację $z=A\setminus B$
wyraża koniunkcja
$\forall u\in z\,(u\in A\land u\notin B)$ oraz
$\forall u\in A\,(u\notin B\Rightarrow u\in z)$.
Dla iloczynu kartezjańskiego sprawdzamy, że każdy jego
element jest parą złożoną z elementów odpowiednich
argumentów i że wszystkie takie pary należą do wyniku.
Kwantyfikatory po parach ograniczamy do argumentów,
a po ich składnikach do sum i zbiorów par: takie
operacje, podobnie jak przestawienia współrzędnych,
mają opisy $\Delta_0$. Dla sumy używamy
$u\in\bigcup A\iff\exists a\in A\;(u\in a)$,
a dla dziedziny
$u\in\operatorname{dom}(A)\iff
\exists p\in A\;\exists v\in\bigcup\bigcup\{p\}\;(p=(u,v))$.
Pozostałe operacje są połączeniami tych samych wzorców.
Indukcja po długości złożenia zachowuje absolutność
każdej operacji otrzymanej w ten sposób.

\emph{Postać normalna.} Dowodzimy przez indukcję po
budowie formuły ograniczonej $\psi$,
że relację
$\{\vec u\in A_1\times\cdots\times A_n:\psi(\vec u)\}$
można otrzymać przez skończone złożenie operacji
$G_i$. Dla formuły atomowej $u_i\in u_j$ używamy
$G_3$ i przestawień współrzędnych; równość sprowadzamy
do dwóch inkluzji z ograniczonymi kwantyfikatorami.
Negacja odpowiada różnicy od całego iloczynu, koniunkcja
-- przecięciu, a kwantyfikator ograniczony
$\exists v\in u_i$ -- rzutowi odpowiedniej relacji
po uprzednim przecięciu z relacją $v\in u_i$.
Rzut realizuje $G_7$ po przestawieniu współrzędnych.
Wszystkie te operacje wymagają skończenie wielu kroków;
nie wprowadzamy kwantyfikatora po nieokreślonej klasie.

Dla dowolnej formuły pierwszego rzędu $\varphi$ jej
relatywizacja do $(A,\in)$ zastępuje wszystkie
kwantyfikatory przez kwantyfikatory ograniczone do $A$.
Otrzymujemy zatem, po dodaniu $A$ oraz parametrów,
wynik skończonego złożenia operacji $G_i$.
To dowodzi inkluzji $\subseteq$.
Dla inkluzji odwrotnej indukcja po długości złożenia
pokazuje, że każda operacja $G_i$ ma definicję
$\Delta_0$ z użytymi argumentami jako parametrami.
Jeśli jej końcowy wynik jest podzbiorem tranzytywnego
$A$, ograniczenie kwantyfikatorów do $A$ daje formułę
definiującą ten podzbiór nad $(A,\in)$.
Dla $A=\varnothing$ obie strony są równe
$\{\varnothing\}$ zgodnie z przyjętą konwencją.
\end{proof}


\begin{lemat}[Jednolity kod $\operatorname{Def}$ klasy $\Sigma_1$]
\label{lem:l-def-sigma1-certificate-75}
W ZF istnieje jedna formuła $\mathsf D(A,B)$ klasy
$\Sigma_1$ taka, że dla każdego \emph{tranzytywnego}
zbioru $A$ i dowolnego $B$ zachodzi
\[
\mathsf D(A,B)\quad\Longleftrightarrow\quad
B=\operatorname{Def}(A).
\]
Dla tranzytywnych argumentów graf funkcji
$A\mapsto\operatorname{Def}(A)$ jest też definiowalny
formułą $\Pi_1$, a zatem jest $\Delta_1$ \emph{nad ZF}.
\end{lemat}
\begin{proof}
Oznaczmy $U_0=A\cup\{A\}$ i zdefiniujmy
jednoznacznie, dla $n\in\omega$,
\[
U_{n+1}=U_n\cup
\{G_i(x,y):x,y\in U_n,\ 1\le i\le10\},
\]
przy czym operacje jednoargumentowe ignorują drugi
argument. Zbiór wszystkich użytych argumentów i wartości
jest objęty jednym $U_{n+1}$. W ZF funkcja
$W:n\mapsto U_n$ istnieje przez rekurencję po $\omega$.
Definiujemy $C=\bigcup_{n\in\omega}U_n$.
Z Lematu~\ref{lem:l-godel-normal-closure} wynika
\[
B=\{z\in C:z\subseteq A\}
  =\operatorname{Def}(A).
\]
Relację ``$W$ jest funkcją o dziedzinie $\omega$,
$W(0)=U_0$ i dla każdego $n\in\omega$ spełnia powyższe
równanie'' można zapisać formułą $\Delta_0$
\emph{z parametrami $W,A,\omega$}. W istocie
$z\in U_{n+1}$ oznacza: $z\in U_n$ albo
$z=G_i(x,y)$ dla pewnych $x,y\in U_n$ i skończenie
wielu indeksów $i$; wszystkie kwantyfikatory są
ograniczone do $U_n$, $U_{n+1}$ lub zbioru indeksów.
Własności $C=\bigcup\operatorname{ran}(W)$ oraz
$B=\{z\in C:z\subseteq A\}$ mają również zapisy
$\Delta_0$ (oba kierunki inkluzji sprawdza się
kwantyfikatorami po $C$ i $B$).
Predykat ``$w=\omega$'' można zapisać formułą
$\Delta_0$, charakteryzując $w$ jako pierwszy
niezerowy graniczny ordynał. Wobec tego zdanie
``istnieją $w,W,C$ spełniające powyższe warunki''
ma postać $\exists w\,\exists W\,\exists C\,\delta$
z $\delta$ klasy $\Delta_0$, czyli jest $\Sigma_1$.
Jednoznaczność $B$ wynika z jednoznaczności
$U_n$ i równości $B=C\cap\mathcal P(A)$.

Aby otrzymać równoważny opis $\Pi_1$,
zakładamy tranzytywność $A$ (warunek $\Delta_0$)
i zapisujemy, że \emph{każdy} zestaw świadków
$w,W,C,B'$ spełniający ograniczone warunki
musi mieć $B'=B$. Jest to formuła $\Pi_1$.
W ZF świadkowie zawsze istnieją, zatem warunek
powszechny jest równoważny poprzedniemu \emph{egzystencjalnemu}. Nie wynika z tego, że ich świadkowie
należą do dowolnego tranzytywnego podzbioru uniwersum.
\end{proof}

\begin{lemat}[Definiowalność $\Delta_1$ całej hierarchii]
\label{lem:l-hierarchy-delta1-derived-75}
W ZF istnieją ustalone formuły
$\Phi_{\Sigma}(\alpha,y)$ klasy $\Sigma_1$ i
$\Phi_{\Pi}(\alpha,y)$ klasy $\Pi_1$ takie, że
ZF dowodzi dla każdej liczby porządkowej $\alpha$:
\[
\Phi_{\Sigma}(\alpha,y)
\ \Longleftrightarrow\
\Phi_{\Pi}(\alpha,y)
\ \Longleftrightarrow\ y=L_\alpha.
\]
Ponadto dla każdego $\alpha$ istnieje dokładnie
jedna funkcja $h_\alpha$ o dziedzinie $\alpha+1$,
która koduje wszystkie etapy konstrukcji do $\alpha$.
\end{lemat}
\begin{proof}
Wybieramy \emph{jeden} predykat $\mathsf D$ z
Lematu~\ref{lem:l-def-sigma1-certificate-75}.
Dla funkcji $h$ o dziedzinie $\alpha+1$
sprawdzamy ograniczonymi kwantyfikatorami warunki:
$h(0)=\varnothing$, na każdym etapie granicznym
$h(\lambda)=\bigcup_{\xi<\lambda}h(\xi)$,
a na etapach następnikowych
$\mathsf D(h(\xi),h(\xi+1))$.
Każde użycie $\mathsf D$ ma świadków $w_\xi,W_\xi,C_\xi$
z poprzedniego lematu.
Ponieważ $\xi$ przebiega \emph{zbiór} $\alpha$,
a świadkowie są wyznaczeni przez kanoniczną rekurencję
po $\omega$, schemat zastępowania w ZF pozwala
zebrać ich rodziny w jedną funkcję $Q$ o dziedzinie
zbioru indeksów następnikowych $\xi<\alpha$.
Jest to zastosowanie zastępowania do
jednoznacznie określonych danych, nie aksjomatu wyboru.

Po włączeniu $Q$ jako dodatkowego świadka wszystkie
warunki dotyczące $h,Q$ można wypisać jako jedną
formułę $\Delta_0$: kwantyfikatory przebiegają
$\alpha$, $h$, $Q$, ich wartości lub zbiory kodujące
operacje, a równości etapów sprawdza się oboma
kierunkami inkluzji. Stąd zdanie
``istnieją $h,Q$ będące poprawnym kodem etapów
z $h(\alpha)=y$'' jest $\Sigma_1$.
Istnienie takich $h,Q$ dla wszystkich ordynałów
wynika z rekurencji pozaskończonej w ZF; jednoznaczność
$h$ wynika z indukcji po $\alpha$ i jednoznaczności
$\operatorname{Def}$. Dlatego równoważna w ZF jest
formuła powszechna: dla \emph{każdych} świadków
$h,Q$ spełniających te same warunki wartość
$h(\alpha)$ wynosi $y$. Ma ona postać $\Pi_1$.
Obie formuły definiują więc graf funkcji
$\alpha\mapsto L_\alpha$ nad ZF.
\end{proof}

\begin{przyklad}[Dlaczego $\Delta_1$ nad ZF nie wystarcza]
\label{ex:l-delta1-not-arbitrary-level-75}
Zbiór $L_\omega=V_\omega$ jest tranzytywny, lecz
nie zawiera $\omega$. Świadek $W$ w formule
$\mathsf D$ powyżej musi być funkcją o dziedzinie
$\omega$; zatem $W\notin L_\omega$, nawet jeśli
$A=\varnothing$ i $\operatorname{Def}(A)$ jest
skończony oraz należy do $L_\omega$.
W szczególności egzystencjalne kodowanie
$\Sigma_1$ \emph{nie jest automatycznie prawdziwe}
w każdym granicznym $L_\theta$.
Równoważność opisów $\Sigma_1$ i $\Pi_1$ jest
twierdzeniem \emph{ZF}, nie twierdzeniem o każdej
tranzytywnej strukturze niespełniającej ZF.
To jest konkretna przeszkoda, którą musi rozwiązać
skończony pakiet kondensacji.
\end{przyklad}

\begin{przyklad}[Absolutność operacji nie daje domknięcia]
\label{ex:l-delta0-not-closure-2026}
Niech $N=\{\varnothing\}$. Jest to niepusty zbiór
tranzytywny. Zdanie $z=G_1(a,b)$ jest $\Delta_0$
— przy argumentach i wartości należących do $N$
jego prawdziwość jest absolutna. Mimo to dla
$a=b=\varnothing\in N$ zachodzi
$G_1(a,b)=\{\varnothing\}\notin N$.
Absolutność \emph{grafu} operacji nie oznacza więc
istnienia jej \emph{wartości} w modelu.
Analogicznie sprawdzenie jednej skończonej tablicy
operacji nie dowodzi, że zbiór wszystkich potrzebnych
tablic jest elementem rozpatrywanego modelu.
\end{przyklad}

\begin{lemat}[Identyfikacja poziomu na podstawie poprawnych kodów]
\label{lem:l-external-codes-identify-2026}
Niech $N$ będzie niepustym zbiorem tranzytywnym,
$\beta=N\cap\operatorname{Ord}$ niezerową graniczną
liczbą porządkową i niech zachodzą dwa warunki:
\begin{enumerate}[label=(\roman*)]
\item dla każdego $\alpha<\beta$ rzeczywista,
  \emph{zewnętrznie obliczona} funkcja
  $h_\alpha=\langle L_\xi:\xi\leq\alpha\rangle$
  należy do $N$;
\item dla każdego $x\in N$ istnieje
  $\alpha<\beta$ takie, że $x\in L_\alpha$.
\end{enumerate}
Wówczas $N=L_\beta$.
\end{lemat}
\begin{proof}
Ustalmy $\alpha<\beta$. Na mocy (i) mamy
$h_\alpha\in N$, a ponieważ $N$ jest tranzytywny,
każda para uporządkowana należąca do $h_\alpha$
i każdy jej składnik należy do $N$. W szczególności
$L_\alpha\in N$, a z tranzytywności otrzymujemy
$L_\alpha\subseteq N$; po zsumowaniu tych
inkluzji po $\alpha<\beta$ mamy
$L_\beta\subseteq N$.
Z kolei (ii) daje każdemu $x\in N$ poziom
$L_\alpha$, do którego należy. Stąd
$N\subseteq\bigcup_{\alpha<\beta}L_\alpha=L_\beta$.
Obie inkluzje kończą dowód. W szczególności nie
utożsamialiśmy tu obiektów \emph{wewnętrznie
zadeklarowanych} jako kody z prawdziwymi kodami:
poprawność została jawnie założona w (i).
\end{proof}

\begin{definicja}[Adekwatność dla konstrukcji $L$]
\label{def:l-adequacy-2026}
Niech $N$ będzie zbiorem tranzytywnym.
Nazywamy go \emph{adekwatnym}, jeżeli:
\begin{enumerate}[label=(A\arabic*)]
\item $N$ jest domknięty na dziesięć operacji
  Gödla z Lematu~\ref{lem:l-godel-normal-closure};
\item dla każdego $U\in N$ zbiór wszystkich
  wartości $G_i(a,b)$ dla $a,b\in U$ i
  $1\leq i\leq10$ (z powtórzeniem argumentu
  dla operacji jednoargumentowych) należy do $N$;
\item dla każdej liczby porządkowej $\alpha\in N$
  \emph{zewnętrznie poprawny} początkowy odcinek
  $\langle L_\xi:\xi<\alpha\rangle$ należy do $N$.
\end{enumerate}
Warunek (A3) jest tu najpierw sformułowany
zewnętrznie, aby nie przemycić jego absolutności
przez samą definicję. Jednolitą formułę
$\mathsf{Hist}$ i pełny dowód (A3) podajemy w
Lemacie~\ref{lem:l-all-stage-histories-78}.
\end{definicja}

\begin{lemat}[Jednolity koszt operacji Gödla]
\label{lem:l-godel-finite-stage-bound-77}
Istnieje \emph{jedna} liczba naturalna $c>0$ taka,
że dla dowolnej liczby porządkowej $\eta$ i każdego
$i\in\{1,\ldots,10\}$, jeżeli wszystkie argumenty
operacji $G_i$ należą do $L_\eta$, to
\[
G_i(\vec a)\in L_{\eta+c}.
\]
Dla operacji jednoargumentowych nie wprowadza się
fikcyjnego dodatkowego argumentu. Dowód nie zakłada
aksjomatu wyboru ani istnienia kodu wszystkich poziomów
$L_\xi$ w jednym ustalonym modelu $L_\theta$.
\end{lemat}
\begin{proof}
Używamy własności udowodnionych przy definiowaniu
hierarchii: każdy $L_\eta$ jest tranzytywny,
$L_\eta\subseteq L_{\eta+1}$, a każdy podzbiór
$L_\rho$ definiowalny nad $(L_\rho,\in)$
z parametrami z $L_\rho$ należy do $L_{\rho+1}$.

Niech wszystkie argumenty leżą w $L_\eta$.
Każdy element argumentu, jak również każdy element
jego elementu, należy do $L_\eta$, dzięki
tranzytywności. Przy kodowaniu Kuratowskiego pary
$\langle u,v\rangle=\{\{u\},\{u,v\}\}$
dla $u,v\in L_\eta$ mamy
$\{u\},\{u,v\}\in L_{\eta+1}$, a więc
$\langle u,v\rangle\in L_{\eta+2}$.
Dla ustalonego kodowania trójek przez zagnieżdżone
pary otrzymujemy analogicznie
$\langle u,v,w\rangle\in L_{\eta+4}$.

Sprawdźmy teraz dziesięć typów operacji.
Elementami $G_1(a,b)$ są $a,b$;
elementami $G_2(a,b)$ i $G_3(a,b)$ są odpowiednie
pary o współrzędnych z argumentów;
elementy $G_4(a,b)$, $G_5(a,b)$ i $G_6(a)$
leżą w $L_\eta$; elementy $G_7(a)$ są
pierwszymi współrzędnymi par należących do $a$.
Elementy $G_8(a)$ są odwróconymi parami, a
$G_9(a)$ i $G_{10}(a)$ są przestawionymi
trójkami o współrzędnych z $L_\eta$.
W każdym z tych przypadków wszystkie elementy
wyniku należą do $L_{\eta+4}$.

Ponadto każda relacja $z\in G_i(\vec a)$ jest
wyrażalna \emph{jedną ustaloną} formułą $\Delta_0$
(Lemat~\ref{lem:l-godel-normal-closure}); jej
prawdziwość z parametrami należącymi do
$L_{\eta+4}$ jest absolutna między tym
tranzytywnym zbiorem a uniwersum.
Zatem $G_i(\vec a)$ jest dokładnie podzbiorem
$L_{\eta+4}$ definiowanym nad
$(L_{\eta+4},\in)$ przez wspomnianą formułę.
Wynika stąd
$G_i(\vec a)\in L_{\eta+5}$.
Można więc przyjąć wspólną stałą $c=6$,
niezależną od $i$, $\eta$ i argumentów.
\end{proof}

\begin{lemat}[Dwa autonomiczne warunki adekwatności]
\label{lem:l-limit-godel-adequacy-77}
Dla każdej niezerowej granicznej liczby porządkowej
$\theta$ zbiór $L_\theta$ spełnia warunki
\emph{(A1)} i \emph{(A2)}
Definicji~\ref{def:l-adequacy-2026}.
Dokładniej, dla każdego $U\in L_\theta$ zbiór
\[
T(U)=\{G_i(a,b):a,b\in U,\ 1\le i\le10\}
\]
(z umową dotyczącą operacji jednoargumentowych)
należy do $L_\theta$.
Sama ta konkluzja nie stwierdza warunku (A3).
\end{lemat}
\begin{proof}
Najpierw weźmy dowolne argumenty $a,b\in L_\theta$.
Istnieje $\eta<\theta$ takie, że $a,b\in L_\eta$:
wybieramy wspólny etap dla skończonej rodziny
argumentów, a następnie, jeżeli trzeba,
zwiększamy go o jeden. Lemat~
\ref{lem:l-godel-finite-stage-bound-77}
daje $G_i(\vec a)\in L_{\eta+6}\subseteq L_\theta$,
gdyż dla granicznej $\theta$ i $\eta<\theta$
zachodzi $\eta+6<\theta$. To dowodzi (A1).

Dla (A2) ustalmy $U\in L_\theta$ i wybierzmy
$\eta<\theta$ z $U\in L_\eta$. Tranzytywność
$L_\eta$ implikuje $U\subseteq L_\eta$.
Stąd, na mocy poprzedniego lematu,
każda wartość $G_i(a,b)$ dla $a,b\in U$
należy do $L_{\eta+6}$, czyli
$T(U)\subseteq L_{\eta+6}$.

Zbiór $T(U)$ jest definiowalny nad
$(L_{\eta+6},\in)$ z parametrem $U$ za pomocą
jednej formuły:
\[
 z\in T(U)\ \longleftrightarrow\
 \bigvee_{i=1}^{10}\ \exists a,b\in U\
                 \bigl(z=G_i(a,b)\bigr).
\]
Każdy graf $G_i$ zastępujemy przy tym jego
formułą $\Delta_0$, a wszystkie kwantyfikatory
po argumentach są ograniczone do $U$.
Ze względu na wcześniejszą inkluzję oraz
absolutność tych grafów formuła definiuje
\emph{dokładnie} $T(U)$, nie tylko jego fragment.
Wobec tego
$T(U)\in L_{\eta+7}\subseteq L_\theta$.
Nie korzystaliśmy z zastępowania wewnątrz
$L_\theta$: zbiór wartości został otrzymany
jako definiowalny podzbiór wcześniejszego poziomu.
\end{proof}

\begin{lemat}[Warunek historii: indeksy skończone i $\omega$]
\label{lem:l-initial-history-omega-77}
Niech $\theta>0$ będzie graniczną liczbą porządkową.
Dla każdej $n<\omega$ rzeczywisty zbiór
$H_n=\langle L_k:k<n\rangle$ należy do $L_\theta$.
Jeżeli $\omega\in L_\theta$, to dodatkowo
$H_\omega=\langle L_n:n<\omega\rangle\in L_\theta$.
\end{lemat}
\begin{proof}
Dla $n<\omega$ wszystkie $L_k=V_k$ są zbiorami
dziedzicznie skończonymi; stąd skończona funkcja
$H_n$ należy do $L_\omega=V_\omega\subseteq L_\theta$.

Dla drugiej części zauważmy najpierw, że
$H_\omega\subseteq L_\omega$:
każda para $\langle n,V_n\rangle$ jest
zbiorem dziedzicznie skończonym.
Pokażemy, że jest to podzbiór
\emph{definiowalny} nad $(L_\omega,\in)$.
Dla $n\in\omega$ formuła $\rho(n,y)$ stwierdza
istnienie skończonej funkcji $f$ o dziedzinie
$n+1$, dla której $f(0)=\varnothing$,
$f(k+1)=\mathcal P(f(k))$ dla każdego $k<n$
oraz $f(n)=y$. Wystąpienie $\mathcal P$ jest
tu skrótem jednej formuły pierwszego rzędu
\[
\forall z\,\bigl(z\in f(k+1)\longleftrightarrow
                 z\subseteq f(k)\bigr).
\]
Ponieważ $f(k)$ jest skończony, każdy jego
podzbiór należy do $L_\omega$; wszystkie
funkcje $f$ potrzebne dla skończonego $n$
również należą do $L_\omega$.
Z tego wynika, że $\rho(n,y)$ jest prawdziwe
w $(L_\omega,\in)$ dokładnie dla $y=V_n$.
Zatem nad $L_\omega$ formuła
\[
\exists n,y\,\bigl(\rho(n,y)\land
  z=\langle n,y\rangle\bigr)
\]
definiuje podzbiór $H_\omega$ zbioru $L_\omega$.
W rezultacie
$H_\omega\in\operatorname{Def}(L_\omega)
=L_{\omega+1}$.
Jeśli $\omega\in L_\theta$, to $\omega<\theta$
i z graniczności $\theta$ wynika
$\omega+1<\theta$, a więc $H_\omega\in L_\theta$.
\end{proof}

\begin{uwaga}[Co brakowało w audycie LXXVII]
\label{rem:l-adequacy-remaining-history-77}
Dowód warunków (A1) i (A2) sam nie pozwala zebrać
historii na etapie granicznym: z przynależności
$H_\xi\in L_\theta$ dla każdego $\xi<\alpha$ nie wynika,
że $\{H_\xi:\xi<\alpha\}\in L_\theta$.
Następne dwa lematy usuwają tę lukę bez zakładania
schematu zastępowania w $L_\theta$. Zamiast zbierać
wartości funkcji na całej $\alpha$, definiujemy
$H_\alpha$ jako \emph{jeden} podzbiór $L_\alpha$;
kluczowa jest absolutność skończonych kodów operacji.
\end{uwaga}

\begin{lemat}[Absolutność skończonych kodów operacji Gödla]
\label{lem:l-finite-godel-term-absoluteness-78}
Istnieją ustalone formuły pierwszego rzędu
$\mathsf{Eval}_G(t,A,\vec p,z)$ oraz
$\mathsf{Def}_G(A,B)$ o następujących własnościach.
Niech $N$ będzie \emph{niepustym tranzytywnym} zbiorem
zamkniętym na operacje $G_1,\ldots,G_{10}$.
Dla każdego tranzytywnego $A\in N$ i $B\in N$ mamy
\[
 (N,\in)\models\mathsf{Def}_G(A,B)
 \quad\Longleftrightarrow\quad B=\operatorname{Def}(A).
\]
Formuła $\mathsf{Eval}_G$ opisuje wartość dowolnego
\emph{skończonego} termu zbudowanego z $G_i$, stałej
$A$ i skończonej krotki parametrów z $A$; jest ona
absolutna między $(N,\in)$ i uniwersum dla tych
argumentów. Nie zakładamy, że $N$ spełnia ZF.
\end{lemat}
\begin{proof}
Z niepustości i tranzytywności $N$ wynika
$\varnothing\in N$: wybieramy element o najmniejszej
randze i stosujemy tranzytywność. Każdy zbiór
dziedzicznie skończony należy do $N$. Rzeczywiście,
jeżeli wszystkie elementy skończonego zbioru $u$
należą do $N$, to wielokrotne zastosowanie operacji
$G_1(a,b)=\{a,b\}$ daje $u\in N$ (dla $u=\varnothing$
korzystamy z poprzedniego zdania). W szczególności
$N$ zawiera wszystkie \emph{standardowe} kody formuł,
termów i skończonych krotek.

Kod termu wybieramy jako skończony \emph{program acykliczny}
$t=(n,c)$: $n$ jest dodatnią skończoną liczbą
porządkową, a funkcja $c$ o dziedzinie $n$
w każdej pozycji $j<n$ wskazuje jedną z instrukcji:
``wstaw $A$'', ``wstaw parametr numer $k$''
albo ``zastosuj $G_i$ do wyników instrukcji
o numerach mniejszych od $j$''.
Dla $G_6,\ldots,G_{10}$ instrukcja ma jeden
argument, a dla pozostałych dwa. Skończoność
$n$ można wyrazić \emph{bez używania $\omega$ jako
parametru}: $n$ jest liczbą porządkową i ani $n$,
ani żaden niezerowy $\xi\le n$ nie jest graniczny.
Jest to warunek absolutny w rozważanych zbiorach
tranzytywnych zawierających wszystkie zbiory
dziedzicznie skończone.

Formuła $\mathsf{Eval}_G(t,A,\vec p,z)$ mówi, że
istnieje funkcja $v$ o dziedzinie $n$, której wartość
$v(j)$ jest wskazaną stałą, parametrem albo
$G_i(v(j_1),v(j_2))$ zgodnie z instrukcją $c(j)$,
przy czym $j_1,j_2<j$, oraz $v(n-1)=z$.
Warunki dla $v$ są jedną formułą, ponieważ indeks
$i$ przebiega ustalony skończony zbiór i każdy
graf $G_i$ ma ustalony opis $\Delta_0$ na mocy
Lematu~\ref{lem:l-godel-normal-closure}.
Dla skończonej krotki parametrów $\vec p$
kod ten ma tylko skończenie wiele instrukcji;
nie kwantyfikujemy po dowolnej klasie operacji.
Ponieważ $N$ jest domknięty na wszystkie
$G_i$, rzeczywiste wartości wszystkich węzłów oraz
cała ich skończona tablica należą do $N$. Odwrotnie,
każda tablica uznana w $N$ jest poprawna zewnętrznie:
wykorzystujemy absolutność grafów $G_i$ i indukcję
po skończonym drzewie. Otrzymujemy żądaną
absolutność $\mathsf{Eval}_G$.

Formuła $\mathsf{Def}_G(A,B)$ stwierdza, że $A$ jest
tranzytywny, każdy element $B$ jest podzbiorem $A$
i jest wynikiem pewnego takiego termu oraz że
\emph{każda} wartość termu z argumentami w
$A\cup\{A\}$, będąca podzbiorem $A$, należy do $B$.
Kwantyfikacja po kodach termów jest zwykłą
kwantyfikacją pierwszego rzędu; warunek, że kod jest
skończonym poprawnym termem, wyrażamy przez relację
jego skończonej tablicy. Dla dowolnego termu i
krotki parametrów z $A$ ich kody należą do $N$,
a wartość termu również należy do $N$. Wobec tego
oba kierunki warunku definiującego $B$ są absolutne.
Z postaci normalnej Gödla
(Lemat~\ref{lem:l-godel-normal-closure}) wynika, że
zewnętrznie $B$ jest dokładnie $\operatorname{Def}(A)$.
Nie twierdzimy, że $\operatorname{Def}(A)\in N$
\emph{dla każdego} $A\in N$: lemat mówi o poprawności
relacji, jeżeli rozpatrywana wartość $B$ już należy
do $N$. To rozróżnienie będzie ważne w dowodzie
następnikowym.
\end{proof}

\begin{lemat}[Wszystkie historie hierarchii mają skończony koszt]
\label{lem:l-all-stage-histories-78}
Dla każdej liczby porządkowej $\alpha$ istnieje
$n_\alpha\in\omega\setminus\{0\}$ takie, że rzeczywista
historia $H_\alpha=\langle L_\xi:\xi<\alpha\rangle$
należy do $L_{\alpha+n_\alpha}$. Co więcej, jeżeli
$\alpha>0$ jest graniczna, to
\[
 H_\alpha\in L_{\alpha+1}.
\]
W konsekwencji dla \emph{każdej} niezerowej granicznej
liczby porządkowej $\theta$ i każdego $\alpha<\theta$
zachodzi $H_\alpha\in L_\theta$, czyli warunek (A3).
\end{lemat}
\begin{proof}
Prowadzimy indukcję pozaskończoną po $\alpha$.
Dla $\alpha=0$ historia jest funkcją pustą, więc
$H_0=\varnothing\in L_1$.

Jeżeli $\alpha=\beta+1$, to
$H_\alpha=H_\beta\cup\{\langle\beta,L_\beta\rangle\}$.
Na mocy indukcji $H_\beta$ należy do pewnego
$L_{\beta+n_\beta}$, natomiast $\beta$ i $L_\beta$
należą do $L_{\beta+1}$. Utworzenie pary
uporządkowanej, singletonu i sumy jest skończonym
złożeniem operacji Gödla. Lemat~
\ref{lem:l-godel-finite-stage-bound-77}
i jego dowód dają zatem pewne $m<\omega$ takie, że
$H_{\beta+1}\in L_{\beta+n_\beta+m}$.
Wybierając dostatecznie duże skończone
$n_{\beta+1}>0$, otrzymujemy tezę dla następnika.

Niech teraz $\alpha>0$ będzie graniczna i
załóżmy tezę dla wszystkich $\xi<\alpha$.
Dla $\xi<\alpha$ także $\xi+1<\alpha$, a więc
$H_{\xi+1}\in L_{\xi+1+n_{\xi+1}}\subseteq L_\alpha$.
Ta ostatnia inkluzja zachodzi, ponieważ
$\alpha$ jest graniczna i dodanie skończonej liczby
do liczby porządkowej mniejszej niż $\alpha$
pozostawia wynik poniżej $\alpha$.
Podobnie $\langle\xi,L_\xi\rangle\in L_\alpha$:
$\xi,L_\xi\in L_{\xi+1}$ i kodowanie pary wymaga
tylko skończenie wielu dalszych poziomów.
Dlatego $H_\alpha\subseteq L_\alpha$.

Zdefiniujmy jedną formułę $\mathsf{Hist}(h,\delta)$:
``$h$ jest funkcją o dziedzinie $\delta$,
$h(0)=\varnothing$ gdy $\delta>0$,
$h(\eta+1)$ spełnia $\mathsf{Def}_G(h(\eta),h(\eta+1))$
a na etapach granicznych $h(\lambda)$ jest sumą
wcześniejszych wartości''. W definicji wymagamy
tranzytywności wszystkich poziomów $h(\eta)$.
Jest to jedna formuła pierwszego rzędu: na każdym
etapie kwantyfikujemy tylko po indeksach, wartościach
$h$ i sprawdzamy ustaloną formułę $\mathsf{Def}_G$.
Dla dowolnej tranzytywnej struktury domkniętej na
operacje Gödla, w szczególności dla
$M=L_\alpha$, każdy kod $h\in M$ spełniający
tę formułę jest \emph{zewnętrznie poprawny}.
Dowodzimy tego indukcją po $\eta<\delta$:
krok zerowy i graniczny są absolutne, a krok
następnikowy wynika z poprzedniego lematu.
To dowodzi także jednoznaczności kodu.

Ponieważ $L_\alpha$ jest domknięty na operacje Gödla,
formuła
\[
 \exists\xi\,\exists h\;\bigl(
  \operatorname{Ord}(\xi)\land
  \mathsf{Hist}(h,\xi+1)\land
  z=\langle\xi,h(\xi)\rangle\bigr)
\]
rozpatrywana w $(L_\alpha,\in)$ definiuje
\emph{dokładnie} $H_\alpha$.
W istocie liczby porządkowe należące do $L_\alpha$
to dokładnie $\xi<\alpha$. Dla każdego z nich
indukcja dostarcza prawdziwy kod $H_{\xi+1}\in L_\alpha$,
więc żaden wymagany element $H_\alpha$ nie jest
pominięty. Odwrotnie, każdy kod $h$ uznany w
$L_\alpha$ jest poprawny zewnętrznie na mocy
absolutności $\mathsf{Def}_G$, więc nie daje
żadnego fałszywego elementu. Ponieważ
$H_\alpha\subseteq L_\alpha$ i jest definiowalny
nad $(L_\alpha,\in)$, mamy
$H_\alpha\in\operatorname{Def}(L_\alpha)=L_{\alpha+1}$.
To kończy krok graniczny.

Na koniec, gdy $\alpha<\theta$ i $\theta$ jest
niezerową graniczną liczbą porządkową,
$\alpha+n_\alpha<\theta$. Zatem
$H_\alpha\in L_{\alpha+n_\alpha}\subseteq L_\theta$.
Warunek (A3) jest udowodniony bez schematu
zastępowania wewnątrz $L_\theta$.
\end{proof}

\begin{lemat}[Redukcja (A3) do poprawności kodu]
\label{lem:l-adequacy-to-external-codes-2026}
Niech $N$ będzie tranzytywnym zbiorem adekwatnym,
$\beta=N\cap\operatorname{Ord}$ niezerową graniczną
liczbą porządkową i załóżmy, że
\begin{enumerate}[label=(\roman*)]
\item dla każdego $x\in N$ istnieje liczba
  porządkowa $\alpha\in N$, dla której $N$
  stwierdza $x\in L_\alpha$;
\item za pomocą warunków (A1)--(A3) poprawność
  wewnętrznego kodu $\langle L_\xi:\xi<\alpha\rangle$
  implikuje jego zgodność z kodem zewnętrznym,
  a definicja $x\in L_\alpha$ jest absolutna
  dla $N$ i metateorii.
\end{enumerate}
Wtedy $N=L_\beta$.
\end{lemat}
\begin{proof}
Dla każdego $\alpha<\beta$, ponieważ $\beta$ jest
graniczny, mamy $\alpha+1<\beta$. Warunek (A3)
zastosowany do $\alpha+1$ dostarcza w $N$ kodu
$\langle L_\xi:\xi\leq\alpha\rangle$.
Warunek (ii) identyfikuje ten kod z zewnętrznym
obiektem; stąd warunek (i) Lematu
\ref{lem:l-external-codes-identify-2026}.
Dla dowolnego $x\in N$ wybieramy $\alpha$ z (i).
Na mocy (ii) jego wewnętrzne świadectwo
$x\in L_\alpha$ jest prawdziwe zewnętrznie.
Otrzymujemy również drugą hipotezę tego lematu,
a więc $N=L_\beta$.
\end{proof}

\begin{lemat}[Skończony pakiet absolutności dla poziomów $L$]
\label{lem:l-finite-condensation-package}
Istnieje jedna \emph{ustalona skończona} koniunkcja
zdań $\sigma_*$ języka teorii mnogości taka, że:
\begin{enumerate}[label=(\roman*)]
\item dla każdej niezerowej granicznej liczby
  porządkowej $\theta$ zachodzi
  $(L_\theta,\in)\models\sigma_*$;
\item jeżeli $N$ jest niepustym zbiorem tranzytywnym,
  $(N,\in)\models\sigma_*$ i
  $\beta=N\cap\operatorname{Ord}$, to $\beta$ jest
  graniczny i $N=L_\beta$.
\end{enumerate}
\end{lemat}
\begin{proof}
Wykorzystujemy ustalone formuły
$\mathsf{Def}_G$ i $\mathsf{Hist}$ z Lematów
\ref{lem:l-finite-godel-term-absoluteness-78} i
\ref{lem:l-all-stage-histories-78}.
Wszystkie poniższe warunki są \emph{pojedynczymi}
zdaniami pierwszego rzędu, nie schematami:
\begin{enumerate}[label=(\arabic*)]
\item domknięcie na dziesięć operacji Gödla:
 dla wszystkich argumentów istnieją ich wartości,
 opisane dziesięcioma ustalonymi grafami $\Delta_0$;
\item każda liczba porządkowa ma w strukturze
 większą liczbę porządkową;
\item dla każdej liczby porządkowej $\alpha$
 istnieje $h$ spełniające
 $\mathsf{Hist}(h,\alpha+1)$;
\item dla każdego $x$ istnieje liczba porządkowa
 $\alpha$ i $h$ spełniające
 $\mathsf{Hist}(h,\alpha+1)$ oraz $x\in h(\alpha)$.
\end{enumerate}
Koniunkcję tych zdań oznaczamy przez $\sigma_*$.
Można do niej dodać warunek (A2), ponieważ
Lemat~\ref{lem:l-limit-godel-adequacy-77}
dowodzi go dla wszystkich rozpatrywanych poziomów;
nie jest on jednak potrzebny do samej identyfikacji
kolapsu.

Niech $\theta>0$ będzie graniczna. Warunek (1)
wynika z Lematu~\ref{lem:l-limit-godel-adequacy-77}.
Warunek (2) wynika z graniczności $\theta$ i
równości $L_\theta\cap\operatorname{Ord}=\theta$.
Dla $\alpha<\theta$ mamy
$H_{\alpha+1}\in L_\theta$ na mocy
Lematu~\ref{lem:l-all-stage-histories-78}, a
absolutność kodu z tego samego lematu zapewnia
$(L_\theta,\in)\models\mathsf{Hist}(H_{\alpha+1},\alpha+1)$.
Stąd (3). Dla $x\in L_\theta$ wybieramy
$\alpha<\theta$ takie, że $x\in L_\alpha$,
co jest możliwe, ponieważ
$L_\theta=\bigcup_{\xi<\theta}L_\xi$.
Ten sam kod $H_{\alpha+1}$ jest świadkiem (4).
Wykazaliśmy (i).

Niech teraz $N$ spełnia założenia (ii).
Ponieważ $N$ jest niepusty i tranzytywny, zawiera
$\varnothing$, a warunek (2) pokazuje, że
$\beta=N\cap\operatorname{Ord}$ jest niezerową
graniczną liczbą porządkową. Na mocy (1) i
Lematu~\ref{lem:l-finite-godel-term-absoluteness-78}
każdy kod $h\in N$ spełniający w $N$ formułę
$\mathsf{Hist}$ jest \emph{zewnętrznie} prawdziwym
odcinkiem hierarchii $L$; wynika to z indukcji
po dziedzinie kodu. Z (3) otrzymujemy dla każdego
$\alpha<\beta$ rzeczywistą historię
$H_{\alpha+1}\in N$. Ponieważ $N$ jest tranzytywny,
$L_\alpha\in N$, a zatem $L_\alpha\subseteq N$.
Po zsumowaniu po $\alpha<\beta$ mamy
$L_\beta\subseteq N$.
Z kolei dla dowolnego $x\in N$ warunek (4)
dostarcza prawdziwego kodu historii i
$\alpha<\beta$ z $x\in L_\alpha$. Stąd
$N\subseteq L_\beta$. Otrzymujemy $N=L_\beta$.
Cały dowód jest wewnętrzny: opiera się na
skończonym kodowaniu operacji, indukcji
pozaskończonej w metateorii ZF i definicji $L$.
\end{proof}

\begin{uwaga}[Zakres uzyskanego dowodu]
\label{rem:l-finite-package-source}
Skończoność $\sigma_*$ wynika z tego, że
$\mathsf{Eval}_G$, $\mathsf{Def}_G$ i $\mathsf{Hist}$
są \emph{ustalonymi} formułami języka teorii
mnogości. Kwantyfikowanie po wszystkich
skończonych kodach termów nie oznacza użycia
predykatu prawdy dla całego uniwersum.
W przeciwieństwie do starszej wersji rozdziału
warunek historii (A3) jest tu udowodniony
(Lemat~\ref{lem:l-all-stage-histories-78}), a nie
przyjmowany jako wynik zewnętrzny.
Zasadnicza zależność tego dowodu to już
udowodniona w rozdziale postać normalna Gödla
(Lemat~\ref{lem:l-godel-normal-closure}).
\end{uwaga}

\begin{lemat}[Dziedzicznie skończony początek kadłuba]
\label{lem:l-hf-fixed-collapse-76}
Niech $\theta>0$ będzie graniczną liczbą porządkową
i $X\preccurlyeq(L_\theta,\in)$.
Wówczas $L_\omega\subseteq X$, a kolaps Mostowskiego
$\pi:X\to N$ spełnia $\pi(a)=a$ dla każdego
$a\in L_\omega$. W szczególności dla $\theta=\omega$
mamy $X=L_\omega$ i $N=L_\omega$.
\end{lemat}
\begin{proof}
Dla dowolnego dziedzicznie skończonego zbioru $a$
konstruujemy przez indukcję po randze formułę
$\delta_a(x)$ bez parametrów, która jednoznacznie
określa $a$ w każdej tranzytywnej strukturze
zawierającej $a$ i jego elementy. Dla $a=\varnothing$
można wziąć $\forall z\,(z\notin x)$.
Jeżeli $a=\{a_0,\ldots,a_{k-1}\}$, to przyjmujemy
\[
\delta_a(x)\equiv\forall z\,
 \bigl(z\in x\longleftrightarrow
  \bigvee_{i<k}\delta_{a_i}(z)\bigr),
\]
a dla $k=0$ zastępujemy alternatywę formułą
fałszywą. Indukcja uzasadnia jednoznaczność:
w tranzytywnym $L_\theta$ każde $a_i$ jest
jedynym rozwiązaniem $\delta_{a_i}$, więc
ekstensjonalność wyznacza jedyny możliwy zbiór $x=a$.
Ponieważ $L_\omega\subseteq L_\theta$ i
$X\preccurlyeq L_\theta$, jedyny świadek zdania
$\exists!x\,\delta_a(x)$ należy do $X$.
Stąd $L_\omega\subseteq X$.

Przez indukcję po randze $a\in L_\omega$
obliczamy kolaps:
$\pi(a)=\{\pi(b):b\in a\cap X\}$.
Wszystkie elementy $b\in a$ są dziedzicznie
skończone, zatem należą do $X$ i przez indukcję
$\pi(b)=b$. Wobec tego $\pi(a)=a$.
Dla $\theta=\omega$ inkluzja
$L_\omega\subseteq X\subseteq L_\omega$
daje równość; w tym przypadku nie potrzebujemy
jednolitego kodowania poziomów konstruowalnych.
\end{proof}

\begin{lemat}[Klasyczny lemat kondensacji Gödla]
\label{lem:l-condensation}
Niech $\theta>0$ będzie graniczną liczbą porządkową,
a $X\preccurlyeq(L_\theta,\in)$ podstrukturą elementarną.
Wówczas kolaps Mostowskiego $(X,\in)$ jest izomorficzny
z $(L_\beta,\in)$ dla pewnej granicznej liczby
porządkowej $\beta\le\theta$.
\end{lemat}
\begin{proof}
Rozważmy $(X,\in)$ z tezy i jego kolaps Mostowskiego
$\pi:X\to N$. Jest on dobrze ufundowany jako
podrelacja rzeczywistej przynależności na zbiorze
$L_\theta$, a ekstensjonalność wynika z
$X\preccurlyeq L_\theta$: formuła odróżniająca dwa
różne zbiory przez element ma świadka w $X$.
Zatem $N$ jest tranzytywnym zbiorem, a $\pi$ jest
izomorfizmem struktur.

Na mocy Lematu~\ref{lem:l-finite-condensation-package}(i)
struktura $L_\theta$ spełnia jedno ustalone zdanie
$\sigma_*$. Pełna elementarność daje
$X\models\sigma_*$, a zachowanie prawdziwości zdań
przez izomorfizm daje $N\models\sigma_*$.
Lemat~\ref{lem:l-finite-condensation-package}(ii)
dostarcza granicznej liczby porządkowej $\beta$
z $N=L_\beta$.

Aby wykazać $\beta\leq\theta$, przypomnijmy,
że dla liczby porządkowej $\gamma\in X$ wartość
$\pi(\gamma)$ jest typem porządkowym zbioru
$X\cap\gamma$, więc ordynały obrazu kolapsu
tworzą podzbiór $\theta$ w sensie porządkowym.
Dokładniej, funkcja $\pi$ ograniczona do
$X\cap\operatorname{Ord}$ jest ściśle rosnąca,
a jej obraz to $\beta$. Ponieważ
$X\cap\operatorname{Ord}\subseteq\theta$,
niemożliwe jest $\beta>\theta$.
Wynika stąd teza.
\end{proof}
\begin{uwaga}[Niezależna kontrola kroku kolapsu]
\label{rem:lvi-condensation-dependency}
Ogólny Lemat~\ref{lem:l-condensation} jest teraz
konsekwencją \emph{udowodnionego wewnętrznie}
skończonego pakietu z
Lematu~\ref{lem:l-finite-condensation-package}.
Nie wykorzystujemy nieuzasadnionej równości
$\pi[X\cap L_\gamma]=L_{\pi(\gamma)}$;
jednolitość kodu jest zapewniona przez
Lematy~\ref{lem:l-finite-godel-term-absoluteness-78}
i~\ref{lem:l-all-stage-histories-78}.
Alternatywny dowód dla kadłuba w języku
rozszerzonym pozostaje w
Lemacie~\ref{lem:l-expanded-hull-condensation};
nie jest już potrzebny do zamknięcia samej kondensacji.
\end{uwaga}

\begin{twierdzenie}[Kondensacja z jednolitego rozpoznawania poziomów]
\label{thm:l-condensation-from-uniform-recognition-75}
Niech $\theta>0$ będzie graniczną liczbą porządkową.
Załóżmy, że istnieje \emph{jedna ustalona formuła}
$\Lambda(\gamma,y)$ języka teorii mnogości taka, że
\begin{enumerate}[label=(\roman*)]
\item dla wszystkich $\gamma<\theta$ mamy
$L_\gamma\in L_\theta$;
\item dla wszystkich liczb porządkowych
$\gamma<\theta$ i $y\in L_\theta$ zachodzi
\[
(L_\theta,\in)\models\Lambda(\gamma,y)
\quad\Longleftrightarrow\quad y=L_\gamma.
\]
\end{enumerate}
Wówczas dla \emph{każdego}
$X\preccurlyeq(L_\theta,\in)$ jego kolaps Mostowskiego
jest dokładnie $(L_\beta,\in)$ dla pewnej
\emph{granicznej} liczby porządkowej $\beta\le\theta$.
\end{twierdzenie}
\begin{proof}
Niech $\pi:X\to N$ będzie kolapsem i
$\beta=N\cap\operatorname{Ord}$.
Elementarność zapewnia $0\in X$ oraz
$\gamma+1\in X$ dla każdej liczby porządkowej
$\gamma\in X$, ponieważ $0$ i operacja następnika
są jednoznacznie definiowalne w $L_\theta$.
Kolaps zachowuje następnik ordynalny, więc $\beta$
jest niezerową liczbą porządkową bez elementu
największego, a zatem jest graniczny.
Dla każdej liczby porządkowej $\gamma\in X$
prawdziwe jest w $L_\theta$ zdanie
$\exists!y\,\Lambda(\gamma,y)$, którego jedynym
świadkiem jest $L_\gamma$. Z elementarności otrzymujemy
$L_\gamma\in X$.

\emph{Indukcja po liczbach porządkowych z $X$.}
Pokażemy, że
\[
\pi(L_\gamma)=L_{\pi(\gamma)}
\qquad(\gamma\in X\cap\operatorname{Ord}).
\]
Dla $\gamma=0$ mamy $L_0=\varnothing$ i $\pi(0)=0$.
Jeżeli $\gamma=\eta+1$, to $\eta\in X$, bo jest
jednoznacznie definiowalnym poprzednikiem
$\gamma$ w $L_\theta$.
Ponieważ $\pi(\gamma)=\pi(\eta)+1$, hipoteza
indukcyjna i Lemat~\ref{lem:l-collapse-def-compatibility}
dają
\begin{align*}
\pi(L_\gamma)
 &=\pi[X\cap\operatorname{Def}(L_\eta)]\\
 &=\operatorname{Def}(L_{\pi(\eta)})
  =L_{\pi(\gamma)}.
\end{align*}
Pierwsza równość wynika z definicji kolapsu
zastosowanej do $L_\gamma\in X$.

Niech teraz $\lambda\in X$ będzie niezerową
graniczną liczbą porządkową.
Dla każdego $z\in X\cap L_\lambda$ w $L_\theta$
zachodzi formuła
``$\exists\xi<\lambda\,\exists y\,
(\Lambda(\xi,y)\land z\in y)$''.
Elementarność daje świadków $\xi,y\in X$;
na mocy (ii) koniecznie $y=L_\xi$.
Otrzymujemy
\[
\pi(L_\lambda)=
 \bigcup_{\xi\in X\cap\lambda}\pi(L_\xi).
\]
Równość odwrotnej inkluzji wynika z
$L_\xi\subseteq L_\lambda$ dla $\xi<\lambda$.
Ponieważ $\pi(\lambda)=\{\pi(\xi):\xi\in X\cap\lambda\}$
i stosujemy hipotezę indukcyjną, prawa strona
jest $\bigcup_{\zeta<\pi(\lambda)}L_\zeta
=L_{\pi(\lambda)}$.

\emph{Identyfikacja całego kolapsu.}
Każde $z\in L_\theta$ należy do pewnego $L_\gamma$
z $\gamma<\theta$. Zatem w $L_\theta$ jest prawdziwe
zdanie $\forall z\,\exists\gamma\,\exists y
(\operatorname{Ord}(\gamma)\land\Lambda(\gamma,y)
\land z\in y)$.
Dla $z\in X$ świadkowie $\gamma,y$ mogą być
wybrane w $X$. Po kolapsie $\pi(z)\in L_{\pi(\gamma)}$,
a więc $N\subseteq L_\beta$.
Z drugiej strony dla każdego $\zeta<\beta$ istnieje
$\gamma\in X\cap\operatorname{Ord}$ z
$\pi(\gamma)=\zeta$. Wówczas
$L_\zeta=\pi(L_\gamma)\in N$, a tranzytywność
$N$ daje $L_\zeta\subseteq N$.
Sumując po $\zeta<\beta$, uzyskujemy
$L_\beta\subseteq N$. Ostatecznie $N=L_\beta$;
nierówność $\beta\le\theta$ wynika z faktu,
że $\pi$ porządkowo odwzorowuje $X\cap\theta$
na $\beta$.
\end{proof}
\begin{uwaga}[Precyzyjna granica redukcji]
\label{rem:l-recognition-gap-75}
Twierdzenie~\ref{thm:l-condensation-from-uniform-recognition-75}
ma kompletny dowód \emph{warunkowy}; nie używa
pakietu $\sigma_*$ ani tablicy spełniania dla
klasy $L$. Lemat~\ref{lem:l-hierarchy-delta1-derived-75}
udowadnia, że nad \emph{pełną ZF} istnieje jedna
formuła definiująca $\gamma\mapsto L_\gamma$.
Nadal jednak nie wykazaliśmy, że dla \emph{każdego}
granicznego $L_\theta$ ta sama formuła rozpoznaje
własne poziomy poprawnie: brak może dotyczyć
wewnętrznych świadków egzystencjalnej definicji.
Przykład~\ref{ex:l-delta1-not-arbitrary-level-75}
pokazuje, dlaczego samo zaklasyfikowanie formuły
do $\Delta_1^{\mathrm{ZF}}$ nie wystarcza.
To jest teraz \emph{pozostający istotny krok wymagany}
w tym alternatywnym dowodzie kondensacji; nie należy
mylić go z elementarnością przekrojów, która jest
już dowiedziona w Lemacie~\ref{lem:l-elementary-section-75}.
\end{uwaga}

\begin{wniosek}[Samodzielna kondensacja na współkońcowo wielu poziomach]
\label{cor:l-condensation-reflecting-levels-75}
Dla każdej liczby porządkowej $\eta$ istnieje
graniczna liczba porządkowa $\theta>\eta$ taka, że
\emph{dla każdej} podstruktury
$X\preccurlyeq(L_\theta,\in)$ jej kolaps Mostowskiego
jest strukturą $(L_\beta,\in)$ dla pewnej
$\beta\le\theta$. Wniosek jest dowiedziony w ZF
bez używania zewnętrznego pakietu adekwatności
z Lematu~\ref{lem:l-finite-condensation-package}.
\end{wniosek}
\begin{proof}
Niech $\Phi_{\Sigma}(\gamma,y)$ będzie
\emph{jedną ustaloną} formułą z
Lematu~\ref{lem:l-hierarchy-delta1-derived-75}.
Ponieważ klasa $L$ spełnia ZF
(Twierdzenie~\ref{thm:l-model-zf}), wewnątrz $L$
ta formuła definiuje graf funkcji
$\gamma\mapsto(L_\gamma)^L$.
Lemat~\ref{lem:l-hierarchy-internal} dowodzi, że
$(L_\gamma)^L=L_\gamma$ także w metateorii ZF.
Zatem dla każdych $\gamma,y\in L$ z $\gamma$
liczbą porządkową mamy
\[
L\models\Phi_{\Sigma}(\gamma,y)
\quad\Longleftrightarrow\quad y=L_\gamma.
\]
Stosujemy Lemat~\ref{lem:l-finite-reflection} do
\emph{skończonej} rodziny składającej się ze
wszystkich podformuł $\Phi_{\Sigma}$, otrzymując
graniczne $\theta>\eta$ z odbiciem tej formuły.
Dla $\gamma<\theta$ zbiór $L_\gamma$ należy do
$L_{\gamma+1}=\operatorname{Def}(L_\gamma)$,
gdyż jest nad sobą definiowany formułą $u=u$.
Ponieważ $\theta$ jest graniczny, mamy
$\gamma+1<\theta$, zatem $L_\gamma\in L_\theta$.
Dla każdego $y\in L_\theta$ odbicie daje
\[
(L_\theta,\in)\models\Phi_{\Sigma}(\gamma,y)
\iff L\models\Phi_{\Sigma}(\gamma,y)
\iff y=L_\gamma.
\]
Znaleźliśmy więc jedną formułę $\Lambda=\Phi_{\Sigma}$
spełniającą obie hipotezy
Twierdzenia~\ref{thm:l-condensation-from-uniform-recognition-75}.
Jego zastosowanie dowodzi tezy \emph{jednocześnie
we wszystkich} podstrukturach elementarnych
wybranego $L_\theta$.
\end{proof}
\begin{uwaga}[Dlaczego wniosek nie zastępuje lematu ogólnego]
\label{rem:l-reflecting-levels-not-all-75}
Odbicie zapewnia \emph{współkońcowo wiele}
odpowiednich poziomów $L_\theta$, nie wszystkie
poziomy graniczne. Wniosek
\ref{cor:l-condensation-reflecting-levels-75}
nie dowodzi zatem bezpośrednio ogólnego
Lematu~\ref{lem:l-condensation}, którego teza
obejmuje \emph{każde} graniczne $\theta$.
Jest natomiast odrębnym, autonomicznym wynikiem
o kondensacji, użytecznym wszędzie tam, gdzie
można swobodnie wybrać dostatecznie duży poziom
za pomocą odbicia.
\end{uwaga}

\begin{lemat}[Liczność poziomów hierarchii $L$]
\label{lem:l-level-cardinality}
W teorii $\mathrm{ZFC}+V=L$, dla każdej nieskończonej liczby
porządkowej $\xi$ zachodzi $|L_\xi|=|\xi|$.
\end{lemat}
\begin{proof}
W każdym kroku hierarchii mamy $L_\xi\cap\operatorname{Ord}=\xi$,
a więc $\xi\subseteq L_\xi$ i $|\xi|\le |L_\xi|$.
Wystarczy dowieść nierówności odwrotnej. Dla $\xi=\omega$
wszystkie poziomy $L_n$ są skończone, zatem $L_\omega$
jest zbiorem przeliczalnym.

Jeśli $\xi=\eta+1$, gdzie $\eta\ge\omega$, każdy zbiór
$z\in L_{\eta+1}=\operatorname{Def}(L_\eta)$ ma kod postaci
$(e,\vec p)$: $e\in\omega$ koduje formułę, a
$\vec p\in L_\eta^{<\omega}$ jest skończoną krotką parametrów.
Wobec tego istnieje surjekcja z pewnego podzbioru
$\omega\times L_\eta^{<\omega}$ na $L_{\eta+1}$.
W ZFC moc zbioru kodów nie przekracza
$\max(\aleph_0,|L_\eta|)=|\eta|$.
Z założenia indukcyjnego wynika
$|L_{\eta+1}|\le|\eta|=|\eta+1|$.

Jeśli $\xi>\omega$ jest graniczny, to
$L_\xi=\bigcup_{\omega\le\eta<\xi}L_\eta$.
Rozważamy zbiór par $(\eta,z)$, gdzie
$\omega\le\eta<\xi$ i $z\in L_\eta$.
Odwzorowanie $(\eta,z)\mapsto z$ jest surjekcją na $L_\xi$.
Na mocy indukcji zbiór par ma moc co najwyżej
$\sum_{\eta<\xi}\max(\aleph_0,|\eta|)\le
|\xi|\cdot|\xi|=|\xi|$, gdzie użyliśmy aksjomatu
wyboru, dostępnego w $V=L$. Otrzymujemy
$|L_\xi|\le|\xi|$, co kończy indukcję.
\end{proof}

\begin{lemat}[Kadłub elementarny o zadanej mocy]
\label{lem:l-skolem-hull-kappa}
Załóżmy $\mathrm{ZFC}+V=L$. Niech $\kappa$ będzie nieskończoną
liczbą kardynalną, $\theta>0$ liczbą porządkową, a
$A\subseteq L_\theta$ zbiorem mocy nie większej niż $\kappa$.
Wówczas istnieje $X\subseteq L_\theta$ takie, że
\[
 A\subseteq X,\qquad X\preccurlyeq (L_\theta,\in),
 \qquad |X|\le\kappa.
\]
Jeżeli $|A|=\kappa$, to $|X|=\kappa$.
\end{lemat}
\begin{proof}
Dla każdej formuły $\exists y\,\psi(y,\vec x)$
ustalamy funkcję Skolema na krotkach $\vec a\in L_\theta$:
jeżeli istnieje świadek $y\in L_\theta$, wybieramy
$<_{L}$-najmniejszego takiego świadka; w przeciwnym
wypadku wybieramy ustalony element $L_\theta$.
Ponieważ $\theta>0$, struktura ma niepusty zbiór
obiektów; funkcje mają więc wartości w $L_\theta$.
Ponieważ język teorii mnogości jest przeliczalny,
funkcji Skolema jest przeliczalnie wiele. Porządek
$<_{L}$ pozwala wyznaczyć świadków bez dodatkowych
arbitralnych wyborów.

Kładziemy $X_0=A$ (lub dodajemy jeden ustalony element,
aby pracować z niepustą strukturą) i indukcyjnie definiujemy
$X_{n+1}$ jako sumę $X_n$ i zbiorów wartości
wszystkich funkcji Skolema na skończonych krotkach
z $X_n$. Wykorzystując arytmetykę nieskończonych liczb
kardynalnych w ZFC, otrzymujemy $|X_n|\le\kappa$
indukcyjnie dla każdego $n\in\omega$. Stąd
$X=\bigcup_{n\in\omega}X_n$ ma moc co najwyżej $\kappa$.
Każdy prawdziwy w $L_\theta$ warunek egzystencjalny
z parametrami w $X$ ma świadka w pewnym $X_{n+1}$.
Z kryterium Tarskiego--Vaughta (konsekwencji
Twierdzenia~\ref{thm:downward-ls}) wynika elementarność
$X\preccurlyeq L_\theta$. Jeżeli $|A|=\kappa$,
inkluzja $A\subseteq X$ daje nierówność odwrotną.
\end{proof}



\begin{lemat}[Niezmienniczość definiowalności przy izomorfizmie]
\label{lem:l-def-isomorphism-74}
Niech $A,B$ będą zbiorami niepustymi, $E\subseteq A\times A$,
$F\subseteq B\times B$, a $f:(A,E)\to(B,F)$ izomorfizmem
struktur dla jednego symbolu relacyjnego. Dla każdej formuły
pierwszego rzędu $\varphi(v_0,\ldots,v_n)$ i każdego
$\vec a\in A^{n+1}$ zachodzi
\[
 (A,E)\models\varphi(\vec a)
 \quad\Longleftrightarrow\quad
 (B,F)\models\varphi(f(\vec a)).
\]
Jeżeli $\operatorname{Def}_E(A)$ oznacza rodzinę podzbiorów
$A$ definiowalnych w $(A,E)$ z parametrami z $A$, to
\[
 \{f[Z]:Z\in\operatorname{Def}_E(A)\}
       =\operatorname{Def}_F(B).
\]
\end{lemat}
\begin{proof}
Dowodzimy pierwszej równoważności indukcyjnie po budowie
$\varphi$. Dla $v_i=v_j$ wynika ona z iniekcyjności $f$.
Dla $v_i\,E\,v_j$ wynika bezpośrednio z definicji
izomorfizmu. Krok negacji polega na zanegowaniu obu stron
równoważności indukcyjnej, a krok koniunkcji na połączeniu
dwóch takich równoważności. W kroku egzystencjalnym
$(A,E)\models\exists v\,\psi(v,\vec a)$ wtedy i tylko
wtedy, gdy istnieje $c\in A$ spełniające $\psi(c,\vec a)$.
Z założenia indukcyjnego jest to równoważne istnieniu
$f(c)\in B$ spełniającego $\psi(f(c),f(\vec a))$;
surjektywność $f$ daje również implikację odwrotną.
Pozostałe spójniki i kwantyfikator ogólny są skrótami
poprzednich konstrukcji.

Jeżeli $Z=\{a\in A:(A,E)\models\psi(a,\vec p)\}$,
to pierwsza część dowodu daje
$f[Z]=\{b\in B:(B,F)\models\psi(b,f(\vec p))\}$,
więc $f[Z]\in\operatorname{Def}_F(B)$. Stosując ten sam
argument do izomorfizmu odwrotnego $f^{-1}$ otrzymujemy
inkluzję przeciwną. W szczególności nie zastępujemy tu
zbioru $Z$ jego przecięciem z innym uniwersum bez
uzasadnienia; przy przejściu do kolapsu użyjemy jawnie
$\pi(z)=\{\pi(u):u\in z\cap X\}$ dla $z\in X$.
\end{proof}
\begin{uwaga}[Granica absolutności w powyższym lemacie]
\label{rem:l-isomorphism-fixed-vs-uniform-74}
Indukcja dotyczy \emph{każdej ustalonej} formuły
$\varphi$ oraz dwóch już określonych struktur zbiorowych.
Nie wynika z niej bez dodatkowej pracy, że poziom
$L_\theta$ posiada wewnętrzny zbiór kodów wszystkich
prawdziwych formuł $(L_\delta,\in)$ dla \emph{każdego}
$\delta<\theta$; taki wniosek wymaga także jednolitej
konstrukcji i dostępności odpowiednich zbiorów kodów
na danym poziomie. Dlatego Lemat~\ref{lem:l-def-isomorphism-74}
usuwa przeskok w kroku izomorfizmu kadłuba
rozszerzonego, ale nie zastępuje technicznego
pakietu z Lematu~\ref{lem:l-finite-condensation-package}.
\end{uwaga}

\begin{lemat}[Kondensacja dla kadłuba w języku rozszerzonym]
\label{lem:l-expanded-hull-condensation}
Niech $\theta$ będzie niezerową graniczną liczbą porządkową.
W metateorii ZF rozważmy strukturę $\mathfrak A_\theta$
o uniwersum $L_\theta$, której język oprócz przynależności
zawiera następujące trzy relacje (ich interpretacje są
\emph{zewnętrznymi} zbiorami krotek):
\begin{itemize}
\item $R(\delta,z)$ wtedy i tylko wtedy, gdy $\delta<\theta$
  jest liczbą porządkową i $z\in L_\delta$;
\item $S(\delta,e,\vec p)$ wtedy i tylko wtedy, gdy
  $0<\delta<\theta$, $e$ jest kodem formuły
  $\varphi_e$ o liczbie zmiennych wolnych zgodnej z
  długością $\vec p\in L_\delta^{<\omega}$ oraz
  $(L_\delta,\in)\models\varphi_e(\vec p)$;
\item dla $0<\delta<\theta$ relacja
  $D(\delta,z,e,\vec p)$ zachodzi wtedy i tylko wtedy,
  gdy $e$ koduje formułę $\varphi_e(x,\vec y)$,
  $\vec p\in L_\delta^{<\omega}$ i
  $z=\{u\in L_\delta:(L_\delta,\in)\models
  \varphi_e(u,\vec p)\}$. Dla $\delta=0$
  przyjmujemy $D(0,z,e,\vec p)$ wtedy i tylko wtedy,
  gdy $z=\varnothing$, $\vec p$ jest krotką pustą
  i $e$ jest kodem formuły z wyróżnioną zmienną $x$;
  w ten sposób $D$ koduje
  $\operatorname{Def}(L_0)=\{\varnothing\}$
  bez relacji spełniania dla pustej struktury.
\end{itemize}
Skończone krotki kodujemy standardowo jako zbiory.
Ponieważ $\theta$ jest graniczny, każdy taki kod
złożony z elementów $L_\delta$ należy do $L_\theta$.
Relacje $R,S,D$ istnieją jako zbiory w metateorii:
ich argumenty pochodzą z ustalonego zbioru $L_\theta$,
a spełnianie dla struktury o zbiorowym uniwersum
$L_\delta$ jest definiowalne w ZF.

Jeżeli $X\preccurlyeq\mathfrak A_\theta$ i
$\pi:X\to\bar X$ jest kolapsem Mostowskiego
\emph{reduktu} $(X,\in)$, to
$\bar X=L_\beta$ dla pewnej granicznej liczby
porządkowej $\beta\le\theta$. Ponadto dla każdej
liczby porządkowej $\delta\in X$ zachodzi
\[
\pi[X\cap L_\delta]=L_{\pi(\delta)}.
\]
Nie zakładamy, że relacje $R,S,D$ są definiowalne
\emph{wewnątrz} $(L_\theta,\in)$.
\end{lemat}
\begin{proof}
\emph{Krok 1: elementarność przekrojów.}
Niech $\delta\in X$ będzie dodatnią liczbą porządkową.
Z kryterium Tarskiego--Vaughta wykażemy
$(X\cap L_\delta,\in)\preccurlyeq(L_\delta,\in)$.
Dla $\delta=0$ obie strony są puste; nie używamy
pojęcia elementarności struktur o pustym uniwersum.
W przypadku $\delta>0$ obie strony mają niepuste
uniwersa: $\varnothing\in X\cap L_\delta$ dzięki
elementarności $X$ i definowalności zbioru pustego.
Ustalmy formułę $\exists y\,\psi(y,\vec p)$
i parametry $\vec p\in(X\cap L_\delta)^{<\omega}$.
Jeżeli $(L_\delta,\in)\models\exists y\,\psi(y,\vec p)$,
to w rozszerzonej strukturze $\mathfrak A_\theta$
prawdziwa jest formuła
\[
\exists y\bigl(R(\delta,y)\land
 S(\delta,\ulcorner\psi\urcorner,(y,\vec p))\bigr).
\]
Kod $\ulcorner\psi\urcorner$ jest liczbą naturalną,
a każda liczba naturalna należy do $X$:
struktura $L_\theta$ zawiera wszystkie skończone
liczby porządkowe, które można kolejno jednoznacznie
definiować. Skończona krotka parametrów i jej kod
również należą do $X$, gdyż $L_\theta$ jest domknięty
na skończone kodowanie, a $X$ jest elementarny.
Elementarność $X$ \emph{w rozszerzonym języku}
daje świadka $y\in X$ spełniającego $R(\delta,y)$,
a więc $y\in X\cap L_\delta$. Jest to dokładnie
warunek Tarskiego--Vaughta. Nie użyliśmy tutaj
wewnętrznej definicji relacji $S$ w $L_\theta$.

\emph{Krok 2: indukcja po liczbach porządkowych z $X$.}
Niech $\delta\in X\cap\operatorname{Ord}$. Wykażemy
\[
(\ast)_\delta:\quad
\pi[X\cap L_\delta]=L_{\pi(\delta)}.
\]
Dla $\delta=0$ obie strony są puste.
Jeżeli $\delta=1$, to $0=\varnothing\in X$
i $L_1=\{0\}$, więc
$\pi[X\cap L_1]=\{0\}=L_1=L_{\pi(1)}$.
Na etapach następnikowych $\delta=\eta+1$ z
$\eta\geq1$ stosujemy niżej etap indukcyjny
oraz elementarność przekrojów niepustych.
Wówczas $\eta\in X$, ponieważ operacja
poprzednika liczby porządkowej następnikowej jest
jednoznacznie definiowalna w $(L_\theta,\in)$.
Załóżmy $(\ast)_\eta$. Weźmy $z\in X\cap L_{\eta+1}$.
Z definicji $L_{\eta+1}$ istnieją kod formuły $e$
i skończona krotka parametrów $\vec p\in L_\eta^{<\omega}$
takie, że $D(\eta,z,e,\vec p)$. Dzięki elementarności
\emph{rozszerzonej} struktury wybieramy
$e,\vec p\in X$. Dla każdego $u\in X\cap L_\eta$
Krok 1 daje równoważność
\[
(L_\eta,\in)\models\varphi_e(u,\vec p)
\iff
(X\cap L_\eta,\in)\models\varphi_e(u,\vec p).
\]
Na mocy Lematu~\ref{lem:l-def-isomorphism-74}
relacja definiowalności przechodzi przez izomorfizm.
Kolaps jest izomorfizmem $X\cap L_\eta$ na
$L_{\pi(\eta)}$ na mocy $(\ast)_\eta$, zatem
\[
\pi(z)=\{v\in L_{\pi(\eta)}:
(L_{\pi(\eta)},\in)\models
\varphi_e(v,\pi(\vec p))\}\in
L_{\pi(\eta)+1}=L_{\pi(\delta)}.
\]
I odwrotnie: każdy $b\in L_{\pi(\eta)+1}$ ma taki
kod $e$ i parametry $\vec q\in L_{\pi(\eta)}$.
Z $(\ast)_\eta$ otrzymujemy parametry
$\vec p\in X\cap L_\eta$ o obrazie $\vec q$.
Relacja $D$ określa jednoznacznie $z\in L_{\eta+1}$;
ponieważ $\eta,e,\vec p\in X$, elementarność daje
$z\in X$. Poprzedni rachunek pokazuje $\pi(z)=b$.
To dowodzi $(\ast)_\delta$ w kroku następnikowym.

Jeżeli $\delta$ jest graniczny i $z\in X\cap L_\delta$,
to $z\in L_\eta$ dla pewnej liczby porządkowej
$\eta<\delta$. Relacja $R$ należy do języka
$\mathfrak A_\theta$, dlatego elementarność daje
takie $\eta\in X$. Z założeń indukcyjnych wynika
\[
\pi[X\cap L_\delta]
=\bigcup_{\eta\in X\cap\delta}L_{\pi(\eta)}
=L_{\pi(\delta)}.
\]
Ostatnia równość korzysta z faktu, że obraz
$X\cap\delta$ przez $\pi$ jest dokładnie
liczbą porządkową $\pi(\delta)$.

\emph{Krok 3: identyfikacja całego kolapsu.}
Dla dowolnego $z\in X$ istnieje $\delta<\theta$
z $z\in L_\delta$, ponieważ $\theta$ jest graniczny.
Elementarność dla relacji $R$ daje takie $\delta\in X$.
Krok 2 pokazuje, że $\pi(z)\in L_{\pi(\delta)}$.
Niech $\beta=\bar X\cap\operatorname{Ord}$.
Z definicji kolapsu dla liczb porządkowych
$\beta$ jest liczbą porządkową, a jego zbiór
poprzedników ma postać
$\{\pi(\delta):\delta\in X\cap\operatorname{Ord}\}$.
Ponieważ $X$ jest domknięty na następniki liczb
porządkowych, $\beta$ jest graniczny.
Stąd $\bar X\subseteq L_\beta$.
Dla każdego $\varepsilon<\beta$ istnieje
$\delta\in X\cap\operatorname{Ord}$
z $\pi(\delta)=\varepsilon$. Z Kroku 2 wynika
$L_\varepsilon\subseteq\bar X$.
Biorąc sumę po $\varepsilon<\beta$,
otrzymujemy $L_\beta\subseteq\bar X$.
Zatem $\bar X=L_\beta$. Ponieważ liczby porządkowe
w $X$ są mniejsze od $\theta$, otrzymujemy
$\beta\le\theta$.
\end{proof}

\begin{uwaga}[Porównanie dwóch samodzielnych dowodów]
\label{rem:l-expanded-vs-full-condensation}
Lemat~\ref{lem:l-expanded-hull-condensation}
konstruuje szczególny kadłub elementarny
w języku rozszerzonym o $R,S,D$ i daje osobną
ścieżkę dowodu GCH. Ogólny
Lemat~\ref{lem:l-condensation} został natomiast
udowodniony przez skończony pakiet rozpoznający
wszystkie graniczne poziomy $L_\theta$.
Obie drogi nie powinny być ze sobą utożsamiane:
pierwsza wymaga wyboru rozszerzonego języka,
a druga dowodu absolutności kodów historii.
\end{uwaga}

\begin{twierdzenie}[GCH w uniwersum konstruowalnym]
\label{thm:l-gch}
W teorii ZF dowodzimy, że $L\models\mathrm{ZFC}+\mathrm{GCH}$.
Dokładniej, dla każdej nieskończonej liczby kardynalnej
$\kappa$ \emph{obliczonej wewnątrz klasy $L$} zachodzi
\[
 L\models 2^\kappa=\kappa^+,
\]
gdzie $\kappa^+$ oznacza następną liczbę kardynalną
\emph{obliczoną wewnątrz $L$}.
\end{twierdzenie}
\begin{proof}
Rozdzielimy dwie płaszczyzny argumentu. Najpierw
przeprowadzimy rachunek w teorii $\mathrm{ZFC}+V=L$.
Następnie zastosujemy go wewnątrz klasy $L$.
Możliwość wykonania tego drugiego kroku wynika
z Twierdzenia~\ref{thm:l-model-zf},
Twierdzenia~\ref{thm:l-definable-wellorder} oraz
Lematu~\ref{lem:l-hierarchy-internal}.

\emph{Etap 1: ustalenie liczby kardynalnej i podzbioru.}
Pracujemy chwilowo przy założeniu $V=L$. Niech $\kappa$
będzie dowolną nieskończoną liczbą kardynalną,
a $A$ dowolnym podzbiorem $\kappa$. Ponieważ $V=L$,
zbiór $A$ należy do pewnego poziomu $L_\eta$.
Wybieramy graniczną liczbę porządkową $\theta>\eta$
tak dużą, że $A,\kappa\in L_\theta$.
Możemy to uczynić, ponieważ hierarchia $L$ jest
ciągła na etapach granicznych i jej suma obejmuje $V$.

\emph{Etap 2: kadłub obejmujący wszystkie liczby
porządkowe poniżej $\kappa$.}
Bierzemy zbiór parametrów
$D=\kappa\cup\{\kappa,A\}$, gdzie $\kappa$ po prawej
stronie rozumiemy jako \emph{zbiór wszystkich liczb
porządkowych mniejszych niż $\kappa$}.
Ponieważ $\kappa$ jest nieskończoną liczbą kardynalną,
$|D|=\kappa$. Stosujemy konstrukcję kadłuba z
Lematu~\ref{lem:l-skolem-hull-kappa} do
\emph{rozszerzonej} struktury $\mathfrak A_\theta$
z Lematu~\ref{lem:l-expanded-hull-condensation}.
Jej język jest skończony, a więc przeliczalny;
zawarte w nim relacje są zbiorami w metateorii.
Otrzymujemy $X\preccurlyeq\mathfrak A_\theta$ takie, że
\[
 \kappa\subseteq X,\qquad \kappa,A\in X,
 \qquad |X|=\kappa.
\]
Uwaga: warunek $\kappa\in X$ nie zastępuje
silniejszego warunku $\kappa\subseteq X$.
Gdybyśmy wiedzieli tylko, że zbiór $\kappa$ należy
do $X$, kolaps mógłby zmienić poszczególne jego
elementy i argument dotyczący zbioru $A$ zawiódłby.

\emph{Etap 3: kolaps i kontrola jego wysokości.}
Niech $\pi:X\to\bar X$ będzie kolapsem Mostowskiego.
Lemat~\ref{lem:l-expanded-hull-condensation} daje
$\bar X=L_\beta$ dla pewnej granicznej liczby porządkowej
$\beta\le\theta$. Nie korzystamy tu z niezweryfikowanego
kroku o wewnętrznym kodowaniu wszystkich poziomów $L$
w dowolnej podstrukturze $(L_\theta,\in)$. Ponieważ $\pi$ jest bijekcją,
$|L_\beta|=|X|=\kappa$. Zbiór $L_\beta$ zawiera
każdą liczbę porządkową mniejszą niż $\beta$;
stąd $|\beta|\le\kappa$. Z definicji następnej
liczby kardynalnej wynika $\beta<\kappa^+$.

Dla każdej liczby porządkowej $\xi\le\kappa$
dowodzimy indukcyjnie $\pi(\xi)=\xi$.
Istotnie, gdy $\xi\in X$ i wszystkie liczby porządkowe
$\zeta<\xi$ należą do $X$, definicja kolapsu daje
\[
 \pi(\xi)=\{\pi(\zeta):\zeta\in X\cap\xi\}
 =\{\zeta:\zeta<\xi\}=\xi.
\]
Warunki te zachodzą, bo $\kappa\subseteq X$
i $\kappa\in X$. W szczególności
$\pi(\xi)=\xi$ dla każdego $\xi<\kappa$.

\emph{Etap 4: kolaps nie zmienia zbioru $A$.}
Ponieważ $A\subseteq\kappa$, wszystkie elementy $A$
są liczbami porządkowymi mniejszymi niż $\kappa$,
a więc należą do $X$. Mamy
\[
 \pi(A)=\{\pi(\xi):\xi\in A\cap X\}
       =\{\xi:\xi\in A\}=A.
\]
Jednocześnie $A\in X$, zatem $\pi(A)\in L_\beta$.
Otrzymujemy $A\in L_\beta\subseteq L_{\kappa^+}$.
Z dowolności podzbioru $A$ wynika
\[
 \mathcal P(\kappa)\subseteq L_{\kappa^+}
 \qquad\text{przy założeniu }V=L.
\]

\emph{Etap 5: obliczenie mocy.}
Lemat~\ref{lem:l-level-cardinality} daje
$|L_{\kappa^+}|=\kappa^+$. Stąd
$2^\kappa=|\mathcal P(\kappa)|\le\kappa^+$.
Twierdzenie Cantora (Twierdzenie~\ref{thm:zf-cantor})
daje nierówność $2^\kappa>\kappa$. Ponieważ
porządek liczb kardynalnych jest liniowy w ZFC,
zachodzi $2^\kappa\ge\kappa^+$. Wobec tego
$2^\kappa=\kappa^+$.

\emph{Etap 6: przeniesienie rachunku do klasy $L$.}
Dowód Etapów 1--5 był dowodem formalnym z aksjomatów
$\mathrm{ZFC}+V=L$, a nie argumentem wymagającym
istnienia przeliczalnego tranzytywnego modelu.
W metateorii ZF klasa $L$ spełnia ZFC, a jej
własna hierarchia konstruowalna jest tożsama
z hierarchią zewnętrzną na mocy
Lematu~\ref{lem:l-hierarchy-internal}.
Ponadto $L\models V=L$. Relatywizując do $L$
każdy krok dowodu z Etapów 1--5, otrzymujemy
GCH dla wszystkich nieskończonych liczb
kardynalnych obliczonych wewnątrz $L$.
Nie stwierdzamy, że liczby kardynalne następniki
w $L$ i w otaczającym uniwersum ZF muszą być równe.
\end{proof}

\begin{przyklad}[Znaczenie warunku $\kappa\subseteq X$]
\label{ex:l-hull-kappa-subset}
Jeśli podstruktura elementarna $X$ zawiera liczbę
porządkową $\kappa$ tylko jako \emph{element}, nie wynika
z tego, że zawiera wszystkie liczby porządkowe
mniejsze niż $\kappa$. Na przykład przeliczalna
podstruktura elementarna dużego poziomu $L_\theta$
może zawierać $\omega_1$ jako element, lecz nie może
zawierać wszystkich elementów $\omega_1$.
Jej kolaps wysyła wówczas $\omega_1$ na przeliczalną
liczbę porządkową. Dlatego w dowodzie GCH
wprowadziliśmy do kadłuba \emph{każdy} element $\kappa$,
a nie tylko sam zbiór $\kappa$.
\end{przyklad}

\begin{wniosek}[Względna niesprzeczność AC, GCH i CH]
\label{cor:l-relative-consistency}
Jeżeli ZF jest niesprzeczna, to niesprzeczna jest
również ZFC+GCH, a więc także ZFC+CH.
\end{wniosek}
\begin{proof}
Wszystkie aksjomaty ZFC+GCH po relatywizacji
kwantyfikatorów do definiowalnej klasy $L$
są twierdzeniami ZF na mocy powyższych wyników.
Gdyby z ZFC+GCH dało się wyprowadzić sprzeczność,
pewien skończony dowód używałby tylko skończenie
wielu instancji jej schematów aksjomatów.
Relatywizując każdy wiersz tego dowodu do $L$
i zastępując aksjomaty udowodnionymi w ZF
odpowiednikami, otrzymalibyśmy dowód sprzeczności
w ZF. Zatem niesprzeczność ZF implikuje
niesprzeczność ZFC+GCH.
\end{proof}


\begin{uwaga}[Status dowodu niezależności CH po audycie]
\label{rem:ch-gch-independent-of-full-condensation}
Dowód Twierdzenia~\ref{thm:l-gch} korzysta bezpośrednio
z Lematu~\ref{lem:l-expanded-hull-condensation}.
Wariant dla dowolnego $X\preccurlyeq L_\theta$
zredukowano osobno do skończonego pakietu
z Lematu~\ref{lem:l-finite-condensation-package}.
Pakiet został już udowodniony w całości w
Lemacie~\ref{lem:l-finite-condensation-package},
a dowód GCH ma dodatkowo niezależną wersję dla
kadłuba w języku rozszerzonym.
\end{uwaga}

\begin{uwaga}[Audyt: status lematu kondensacji]
\label{rem:l-condensation-proof-scope}
Bezpośredni dowód Lematu~\ref{lem:l-condensation} można oprzeć
na zgodności operacji $\operatorname{Def}$ z kolapsem, lecz
nie wynika ona z samej elementarności podstruktury. Jeżeli $\pi:X\to\bar X$ jest kolapsem, to
$\pi$ zachowuje relację przynależności między elementami $X$;
nie wolno jednak bez dodatkowego uzasadnienia utożsamiać
$\pi[X\cap L_\gamma]$ z $L_{\pi(\gamma)}$.
Przekrój $X\cap L_\gamma$ nie musi być całym $L_\gamma$,
ale przy $L_\gamma\in X$ jest jego podstrukturą elementarną
na mocy Lematu~
\ref{lem:l-elementary-section-75}.
Prawidłowe uzasadnienie kroku następnikowego wykorzystuje
jednolity kod formuły i krotki parametrów, elementarność
$X$ oraz absolutność relacji spełniania dla zbiorowej struktury.
To rozróżnienie jest istotne przy stosowaniu lematu do GCH.
\end{uwaga}



\begin{uwaga}[Od lokalnych tablic do jednolitego kodu]
\label{rem:l-condensation-scope-76}
Lemat~\ref{lem:l-local-satisfaction-bound-76}
kontroluje tablice ustalonych formuł, a
Lemat~\ref{lem:l-hf-fixed-collapse-76}
przypadek $L_\omega$.
Dopiero Lemat~\ref{lem:l-all-stage-histories-78}
zapewnia wszystkie historie jednocześnie w sensie
\emph{jednej formuły} pierwszego rzędu.
Konkretnie można przyjąć
\[
\Lambda(\gamma,y)\;\equiv\;
 \exists h\bigl(\mathsf{Hist}(h,\gamma+1)
                     \land h(\gamma)=y\bigr).
\]
Na mocy absolutności z
Lematu~\ref{lem:l-finite-godel-term-absoluteness-78}
dla każdego granicznego $L_\theta$ i
$\gamma<\theta$ struktura $(L_\theta,\in)$
rozpoznaje przez $\Lambda(\gamma,y)$ dokładnie
$y=L_\gamma$. Świadek $H_{\gamma+1}$ należy
do $L_\theta$ dzięki
Lematowi~\ref{lem:l-all-stage-histories-78}.
Spełnione są więc także hipotezy
Twierdzenia~\ref{thm:l-condensation-from-uniform-recognition-75}.
\end{uwaga}

\section{Jednolite kody poziomów konstruowalnych — zakres formalizacji}
\label{sec:l-uniform-codes-32}
\begin{definicja}[Kod zbiorowego etapu konstrukcji]
\label{def:l-stage-code-32}
Niech $\alpha$ będzie liczbą porządkową. Przyjmujemy
$\operatorname{Def}(\varnothing)=\{\varnothing\}$
zgodnie z Definicją~\ref{def:constructible-def}. \emph{Kodem etapu} nazywamy funkcję $h$ o dziedzinie $\alpha+1$, dla której $h(0)=\varnothing$, $h(\xi+1)=\operatorname{Def}(h(\xi))$ dla $\xi<\alpha$ i $h(\lambda)=\bigcup_{\xi<\lambda}h(\xi)$ dla każdej niezerowej granicznej liczby porządkowej $\lambda\leq\alpha$. Dla niepustego $A$ relacja $\operatorname{Def}(A)=B$
oznacza, że $B$ jest zbiorem wszystkich podzbiorów $A$
definiowalnych w strukturze $(A,\in)$ przez formuły
pierwszego rzędu z parametrami z $A$; przypadek pusty
reguluje wskazana konwencja.
\end{definicja}
\begin{lemat}[Jednoznaczność kodu etapu]
\label{lem:l-stage-code-unique-32}
W ZF dla każdej liczby porządkowej $\alpha$ istnieje dokładnie jeden kod etapu $h$ o dziedzinie $\alpha+1$; jego ostatnia wartość wynosi $L_\alpha$.
\end{lemat}
\begin{proof}
Dla zbioru $A$ relacja spełniania formuł w zbiorowej strukturze $(A,\in)$ jest definiowalna za pomocą rekurencji po skończonej budowie formuł. Zatem $\operatorname{Def}(A)$ istnieje na mocy aksjomatów zbioru potęgowego i wyodrębniania: wybieramy spośród podzbiorów $A$ te, dla których istnieje kod formuły i skończona krotka parametrów z $A$ definiujące dokładnie dany podzbiór. Jest to operacja jednoznaczna. Twierdzenie o rekurencji pozaskończonej, stosowane do tej operacji i do sum na etapach granicznych, daje funkcję $h$ o wskazanej dziedzinie. Jeśli $h$ i $k$ są dwoma takimi funkcjami, to indukcja pozaskończona po $\xi\leq\alpha$ dowodzi $h(\xi)=k(\xi)$: dla zera wynika to z definicji, dla następnika z jednoznaczności $\operatorname{Def}$, a dla etapu granicznego z równości wcześniejszych wartości. Wartość końcowa jest więc z definicji $L_\alpha$.
\end{proof}
\begin{uwaga}[Formalizacja i jej niezależna kontrola]
\label{rem:l-stage-internality-open-32}
Istnienie kodów z poprzedniego lematu wynika
z rekurencji w metateorii ZF. Dodatkowa treść,
niezbędna do kondensacji, została udowodniona
w Lematach~\ref{lem:l-finite-godel-term-absoluteness-78}
i~\ref{lem:l-all-stage-histories-78}:
każdy graniczny $L_\theta$ zawiera właściwe
kody etapów oraz prawidłowo interpretuje
jedną formułę $\mathsf{Hist}$. Pozwala to używać
skończonego pakietu z
Lematu~\ref{lem:l-finite-condensation-package}
bez założenia ZF wewnątrz $L_\theta$.
\end{uwaga}